\documentclass[11pt]{article}

    \usepackage[breakable]{tcolorbox}
    \usepackage{parskip} % Stop auto-indenting (to mimic markdown behaviour)
    

    % Basic figure setup, for now with no caption control since it's done
    % automatically by Pandoc (which extracts ![](path) syntax from Markdown).
    \usepackage{graphicx}
    % Keep aspect ratio if custom image width or height is specified
    \setkeys{Gin}{keepaspectratio}
    % Maintain compatibility with old templates. Remove in nbconvert 6.0
    \let\Oldincludegraphics\includegraphics
    % Ensure that by default, figures have no caption (until we provide a
    % proper Figure object with a Caption API and a way to capture that
    % in the conversion process - todo).
    \usepackage{caption}
    \DeclareCaptionFormat{nocaption}{}
    \captionsetup{format=nocaption,aboveskip=0pt,belowskip=0pt}

    \usepackage{float}
    \floatplacement{figure}{H} % forces figures to be placed at the correct location
    \usepackage{xcolor} % Allow colors to be defined
    \usepackage{enumerate} % Needed for markdown enumerations to work
    \usepackage{geometry} % Used to adjust the document margins
    \usepackage{amsmath} % Equations
    \usepackage{amssymb} % Equations
    \usepackage{textcomp} % defines textquotesingle
    % Hack from http://tex.stackexchange.com/a/47451/13684:
    \AtBeginDocument{%
        \def\PYZsq{\textquotesingle}% Upright quotes in Pygmentized code
    }
    \usepackage{upquote} % Upright quotes for verbatim code
    \usepackage{eurosym} % defines \euro

    \usepackage{iftex}
    \ifPDFTeX
        \usepackage[T1]{fontenc}
        \IfFileExists{alphabeta.sty}{
              \usepackage{alphabeta}
          }{
              \usepackage[mathletters]{ucs}
              \usepackage[utf8x]{inputenc}
          }
    \else
        \usepackage{fontspec}
        \usepackage{unicode-math}
    \fi

    \usepackage{fancyvrb} % verbatim replacement that allows latex
    \usepackage{grffile} % extends the file name processing of package graphics
                         % to support a larger range
    \makeatletter % fix for old versions of grffile with XeLaTeX
    \@ifpackagelater{grffile}{2019/11/01}
    {
      % Do nothing on new versions
    }
    {
      \def\Gread@@xetex#1{%
        \IfFileExists{"\Gin@base".bb}%
        {\Gread@eps{\Gin@base.bb}}%
        {\Gread@@xetex@aux#1}%
      }
    }
    \makeatother
    \usepackage[Export]{adjustbox} % Used to constrain images to a maximum size
    \adjustboxset{max size={0.9\linewidth}{0.9\paperheight}}

    % The hyperref package gives us a pdf with properly built
    % internal navigation ('pdf bookmarks' for the table of contents,
    % internal cross-reference links, web links for URLs, etc.)
    \usepackage{hyperref}
    % The default LaTeX title has an obnoxious amount of whitespace. By default,
    % titling removes some of it. It also provides customization options.
    \usepackage{titling}
    \usepackage{longtable} % longtable support required by pandoc >1.10
    \usepackage{booktabs}  % table support for pandoc > 1.12.2
    \usepackage{array}     % table support for pandoc >= 2.11.3
    \usepackage{calc}      % table minipage width calculation for pandoc >= 2.11.1
    \usepackage[inline]{enumitem} % IRkernel/repr support (it uses the enumerate* environment)
    \usepackage[normalem]{ulem} % ulem is needed to support strikethroughs (\sout)
                                % normalem makes italics be italics, not underlines
    \usepackage{soul}      % strikethrough (\st) support for pandoc >= 3.0.0
    \usepackage{mathrsfs}
    

    
    % Colors for the hyperref package
    \definecolor{urlcolor}{rgb}{0,.145,.698}
    \definecolor{linkcolor}{rgb}{.71,0.21,0.01}
    \definecolor{citecolor}{rgb}{.12,.54,.11}

    % ANSI colors
    \definecolor{ansi-black}{HTML}{3E424D}
    \definecolor{ansi-black-intense}{HTML}{282C36}
    \definecolor{ansi-red}{HTML}{E75C58}
    \definecolor{ansi-red-intense}{HTML}{B22B31}
    \definecolor{ansi-green}{HTML}{00A250}
    \definecolor{ansi-green-intense}{HTML}{007427}
    \definecolor{ansi-yellow}{HTML}{DDB62B}
    \definecolor{ansi-yellow-intense}{HTML}{B27D12}
    \definecolor{ansi-blue}{HTML}{208FFB}
    \definecolor{ansi-blue-intense}{HTML}{0065CA}
    \definecolor{ansi-magenta}{HTML}{D160C4}
    \definecolor{ansi-magenta-intense}{HTML}{A03196}
    \definecolor{ansi-cyan}{HTML}{60C6C8}
    \definecolor{ansi-cyan-intense}{HTML}{258F8F}
    \definecolor{ansi-white}{HTML}{C5C1B4}
    \definecolor{ansi-white-intense}{HTML}{A1A6B2}
    \definecolor{ansi-default-inverse-fg}{HTML}{FFFFFF}
    \definecolor{ansi-default-inverse-bg}{HTML}{000000}

    % common color for the border for error outputs.
    \definecolor{outerrorbackground}{HTML}{FFDFDF}

    % commands and environments needed by pandoc snippets
    % extracted from the output of `pandoc -s`
    \providecommand{\tightlist}{%
      \setlength{\itemsep}{0pt}\setlength{\parskip}{0pt}}
    \DefineVerbatimEnvironment{Highlighting}{Verbatim}{commandchars=\\\{\}}
    % Add ',fontsize=\small' for more characters per line
    \newenvironment{Shaded}{}{}
    \newcommand{\KeywordTok}[1]{\textcolor[rgb]{0.00,0.44,0.13}{\textbf{{#1}}}}
    \newcommand{\DataTypeTok}[1]{\textcolor[rgb]{0.56,0.13,0.00}{{#1}}}
    \newcommand{\DecValTok}[1]{\textcolor[rgb]{0.25,0.63,0.44}{{#1}}}
    \newcommand{\BaseNTok}[1]{\textcolor[rgb]{0.25,0.63,0.44}{{#1}}}
    \newcommand{\FloatTok}[1]{\textcolor[rgb]{0.25,0.63,0.44}{{#1}}}
    \newcommand{\CharTok}[1]{\textcolor[rgb]{0.25,0.44,0.63}{{#1}}}
    \newcommand{\StringTok}[1]{\textcolor[rgb]{0.25,0.44,0.63}{{#1}}}
    \newcommand{\CommentTok}[1]{\textcolor[rgb]{0.38,0.63,0.69}{\textit{{#1}}}}
    \newcommand{\OtherTok}[1]{\textcolor[rgb]{0.00,0.44,0.13}{{#1}}}
    \newcommand{\AlertTok}[1]{\textcolor[rgb]{1.00,0.00,0.00}{\textbf{{#1}}}}
    \newcommand{\FunctionTok}[1]{\textcolor[rgb]{0.02,0.16,0.49}{{#1}}}
    \newcommand{\RegionMarkerTok}[1]{{#1}}
    \newcommand{\ErrorTok}[1]{\textcolor[rgb]{1.00,0.00,0.00}{\textbf{{#1}}}}
    \newcommand{\NormalTok}[1]{{#1}}

    % Additional commands for more recent versions of Pandoc
    \newcommand{\ConstantTok}[1]{\textcolor[rgb]{0.53,0.00,0.00}{{#1}}}
    \newcommand{\SpecialCharTok}[1]{\textcolor[rgb]{0.25,0.44,0.63}{{#1}}}
    \newcommand{\VerbatimStringTok}[1]{\textcolor[rgb]{0.25,0.44,0.63}{{#1}}}
    \newcommand{\SpecialStringTok}[1]{\textcolor[rgb]{0.73,0.40,0.53}{{#1}}}
    \newcommand{\ImportTok}[1]{{#1}}
    \newcommand{\DocumentationTok}[1]{\textcolor[rgb]{0.73,0.13,0.13}{\textit{{#1}}}}
    \newcommand{\AnnotationTok}[1]{\textcolor[rgb]{0.38,0.63,0.69}{\textbf{\textit{{#1}}}}}
    \newcommand{\CommentVarTok}[1]{\textcolor[rgb]{0.38,0.63,0.69}{\textbf{\textit{{#1}}}}}
    \newcommand{\VariableTok}[1]{\textcolor[rgb]{0.10,0.09,0.49}{{#1}}}
    \newcommand{\ControlFlowTok}[1]{\textcolor[rgb]{0.00,0.44,0.13}{\textbf{{#1}}}}
    \newcommand{\OperatorTok}[1]{\textcolor[rgb]{0.40,0.40,0.40}{{#1}}}
    \newcommand{\BuiltInTok}[1]{{#1}}
    \newcommand{\ExtensionTok}[1]{{#1}}
    \newcommand{\PreprocessorTok}[1]{\textcolor[rgb]{0.74,0.48,0.00}{{#1}}}
    \newcommand{\AttributeTok}[1]{\textcolor[rgb]{0.49,0.56,0.16}{{#1}}}
    \newcommand{\InformationTok}[1]{\textcolor[rgb]{0.38,0.63,0.69}{\textbf{\textit{{#1}}}}}
    \newcommand{\WarningTok}[1]{\textcolor[rgb]{0.38,0.63,0.69}{\textbf{\textit{{#1}}}}}


    % Define a nice break command that doesn't care if a line doesn't already
    % exist.
    \def\br{\hspace*{\fill} \\* }
    % Math Jax compatibility definitions
    \def\gt{>}
    \def\lt{<}
    \let\Oldtex\TeX
    \let\Oldlatex\LaTeX
    \renewcommand{\TeX}{\textrm{\Oldtex}}
    \renewcommand{\LaTeX}{\textrm{\Oldlatex}}
    % Document parameters
    % Document title
    \title{trabalho2}
    
    
    
    
    
    
    
% Pygments definitions
\makeatletter
\def\PY@reset{\let\PY@it=\relax \let\PY@bf=\relax%
    \let\PY@ul=\relax \let\PY@tc=\relax%
    \let\PY@bc=\relax \let\PY@ff=\relax}
\def\PY@tok#1{\csname PY@tok@#1\endcsname}
\def\PY@toks#1+{\ifx\relax#1\empty\else%
    \PY@tok{#1}\expandafter\PY@toks\fi}
\def\PY@do#1{\PY@bc{\PY@tc{\PY@ul{%
    \PY@it{\PY@bf{\PY@ff{#1}}}}}}}
\def\PY#1#2{\PY@reset\PY@toks#1+\relax+\PY@do{#2}}

\@namedef{PY@tok@w}{\def\PY@tc##1{\textcolor[rgb]{0.73,0.73,0.73}{##1}}}
\@namedef{PY@tok@c}{\let\PY@it=\textit\def\PY@tc##1{\textcolor[rgb]{0.24,0.48,0.48}{##1}}}
\@namedef{PY@tok@cp}{\def\PY@tc##1{\textcolor[rgb]{0.61,0.40,0.00}{##1}}}
\@namedef{PY@tok@k}{\let\PY@bf=\textbf\def\PY@tc##1{\textcolor[rgb]{0.00,0.50,0.00}{##1}}}
\@namedef{PY@tok@kp}{\def\PY@tc##1{\textcolor[rgb]{0.00,0.50,0.00}{##1}}}
\@namedef{PY@tok@kt}{\def\PY@tc##1{\textcolor[rgb]{0.69,0.00,0.25}{##1}}}
\@namedef{PY@tok@o}{\def\PY@tc##1{\textcolor[rgb]{0.40,0.40,0.40}{##1}}}
\@namedef{PY@tok@ow}{\let\PY@bf=\textbf\def\PY@tc##1{\textcolor[rgb]{0.67,0.13,1.00}{##1}}}
\@namedef{PY@tok@nb}{\def\PY@tc##1{\textcolor[rgb]{0.00,0.50,0.00}{##1}}}
\@namedef{PY@tok@nf}{\def\PY@tc##1{\textcolor[rgb]{0.00,0.00,1.00}{##1}}}
\@namedef{PY@tok@nc}{\let\PY@bf=\textbf\def\PY@tc##1{\textcolor[rgb]{0.00,0.00,1.00}{##1}}}
\@namedef{PY@tok@nn}{\let\PY@bf=\textbf\def\PY@tc##1{\textcolor[rgb]{0.00,0.00,1.00}{##1}}}
\@namedef{PY@tok@ne}{\let\PY@bf=\textbf\def\PY@tc##1{\textcolor[rgb]{0.80,0.25,0.22}{##1}}}
\@namedef{PY@tok@nv}{\def\PY@tc##1{\textcolor[rgb]{0.10,0.09,0.49}{##1}}}
\@namedef{PY@tok@no}{\def\PY@tc##1{\textcolor[rgb]{0.53,0.00,0.00}{##1}}}
\@namedef{PY@tok@nl}{\def\PY@tc##1{\textcolor[rgb]{0.46,0.46,0.00}{##1}}}
\@namedef{PY@tok@ni}{\let\PY@bf=\textbf\def\PY@tc##1{\textcolor[rgb]{0.44,0.44,0.44}{##1}}}
\@namedef{PY@tok@na}{\def\PY@tc##1{\textcolor[rgb]{0.41,0.47,0.13}{##1}}}
\@namedef{PY@tok@nt}{\let\PY@bf=\textbf\def\PY@tc##1{\textcolor[rgb]{0.00,0.50,0.00}{##1}}}
\@namedef{PY@tok@nd}{\def\PY@tc##1{\textcolor[rgb]{0.67,0.13,1.00}{##1}}}
\@namedef{PY@tok@s}{\def\PY@tc##1{\textcolor[rgb]{0.73,0.13,0.13}{##1}}}
\@namedef{PY@tok@sd}{\let\PY@it=\textit\def\PY@tc##1{\textcolor[rgb]{0.73,0.13,0.13}{##1}}}
\@namedef{PY@tok@si}{\let\PY@bf=\textbf\def\PY@tc##1{\textcolor[rgb]{0.64,0.35,0.47}{##1}}}
\@namedef{PY@tok@se}{\let\PY@bf=\textbf\def\PY@tc##1{\textcolor[rgb]{0.67,0.36,0.12}{##1}}}
\@namedef{PY@tok@sr}{\def\PY@tc##1{\textcolor[rgb]{0.64,0.35,0.47}{##1}}}
\@namedef{PY@tok@ss}{\def\PY@tc##1{\textcolor[rgb]{0.10,0.09,0.49}{##1}}}
\@namedef{PY@tok@sx}{\def\PY@tc##1{\textcolor[rgb]{0.00,0.50,0.00}{##1}}}
\@namedef{PY@tok@m}{\def\PY@tc##1{\textcolor[rgb]{0.40,0.40,0.40}{##1}}}
\@namedef{PY@tok@gh}{\let\PY@bf=\textbf\def\PY@tc##1{\textcolor[rgb]{0.00,0.00,0.50}{##1}}}
\@namedef{PY@tok@gu}{\let\PY@bf=\textbf\def\PY@tc##1{\textcolor[rgb]{0.50,0.00,0.50}{##1}}}
\@namedef{PY@tok@gd}{\def\PY@tc##1{\textcolor[rgb]{0.63,0.00,0.00}{##1}}}
\@namedef{PY@tok@gi}{\def\PY@tc##1{\textcolor[rgb]{0.00,0.52,0.00}{##1}}}
\@namedef{PY@tok@gr}{\def\PY@tc##1{\textcolor[rgb]{0.89,0.00,0.00}{##1}}}
\@namedef{PY@tok@ge}{\let\PY@it=\textit}
\@namedef{PY@tok@gs}{\let\PY@bf=\textbf}
\@namedef{PY@tok@ges}{\let\PY@bf=\textbf\let\PY@it=\textit}
\@namedef{PY@tok@gp}{\let\PY@bf=\textbf\def\PY@tc##1{\textcolor[rgb]{0.00,0.00,0.50}{##1}}}
\@namedef{PY@tok@go}{\def\PY@tc##1{\textcolor[rgb]{0.44,0.44,0.44}{##1}}}
\@namedef{PY@tok@gt}{\def\PY@tc##1{\textcolor[rgb]{0.00,0.27,0.87}{##1}}}
\@namedef{PY@tok@err}{\def\PY@bc##1{{\setlength{\fboxsep}{\string -\fboxrule}\fcolorbox[rgb]{1.00,0.00,0.00}{1,1,1}{\strut ##1}}}}
\@namedef{PY@tok@kc}{\let\PY@bf=\textbf\def\PY@tc##1{\textcolor[rgb]{0.00,0.50,0.00}{##1}}}
\@namedef{PY@tok@kd}{\let\PY@bf=\textbf\def\PY@tc##1{\textcolor[rgb]{0.00,0.50,0.00}{##1}}}
\@namedef{PY@tok@kn}{\let\PY@bf=\textbf\def\PY@tc##1{\textcolor[rgb]{0.00,0.50,0.00}{##1}}}
\@namedef{PY@tok@kr}{\let\PY@bf=\textbf\def\PY@tc##1{\textcolor[rgb]{0.00,0.50,0.00}{##1}}}
\@namedef{PY@tok@bp}{\def\PY@tc##1{\textcolor[rgb]{0.00,0.50,0.00}{##1}}}
\@namedef{PY@tok@fm}{\def\PY@tc##1{\textcolor[rgb]{0.00,0.00,1.00}{##1}}}
\@namedef{PY@tok@vc}{\def\PY@tc##1{\textcolor[rgb]{0.10,0.09,0.49}{##1}}}
\@namedef{PY@tok@vg}{\def\PY@tc##1{\textcolor[rgb]{0.10,0.09,0.49}{##1}}}
\@namedef{PY@tok@vi}{\def\PY@tc##1{\textcolor[rgb]{0.10,0.09,0.49}{##1}}}
\@namedef{PY@tok@vm}{\def\PY@tc##1{\textcolor[rgb]{0.10,0.09,0.49}{##1}}}
\@namedef{PY@tok@sa}{\def\PY@tc##1{\textcolor[rgb]{0.73,0.13,0.13}{##1}}}
\@namedef{PY@tok@sb}{\def\PY@tc##1{\textcolor[rgb]{0.73,0.13,0.13}{##1}}}
\@namedef{PY@tok@sc}{\def\PY@tc##1{\textcolor[rgb]{0.73,0.13,0.13}{##1}}}
\@namedef{PY@tok@dl}{\def\PY@tc##1{\textcolor[rgb]{0.73,0.13,0.13}{##1}}}
\@namedef{PY@tok@s2}{\def\PY@tc##1{\textcolor[rgb]{0.73,0.13,0.13}{##1}}}
\@namedef{PY@tok@sh}{\def\PY@tc##1{\textcolor[rgb]{0.73,0.13,0.13}{##1}}}
\@namedef{PY@tok@s1}{\def\PY@tc##1{\textcolor[rgb]{0.73,0.13,0.13}{##1}}}
\@namedef{PY@tok@mb}{\def\PY@tc##1{\textcolor[rgb]{0.40,0.40,0.40}{##1}}}
\@namedef{PY@tok@mf}{\def\PY@tc##1{\textcolor[rgb]{0.40,0.40,0.40}{##1}}}
\@namedef{PY@tok@mh}{\def\PY@tc##1{\textcolor[rgb]{0.40,0.40,0.40}{##1}}}
\@namedef{PY@tok@mi}{\def\PY@tc##1{\textcolor[rgb]{0.40,0.40,0.40}{##1}}}
\@namedef{PY@tok@il}{\def\PY@tc##1{\textcolor[rgb]{0.40,0.40,0.40}{##1}}}
\@namedef{PY@tok@mo}{\def\PY@tc##1{\textcolor[rgb]{0.40,0.40,0.40}{##1}}}
\@namedef{PY@tok@ch}{\let\PY@it=\textit\def\PY@tc##1{\textcolor[rgb]{0.24,0.48,0.48}{##1}}}
\@namedef{PY@tok@cm}{\let\PY@it=\textit\def\PY@tc##1{\textcolor[rgb]{0.24,0.48,0.48}{##1}}}
\@namedef{PY@tok@cpf}{\let\PY@it=\textit\def\PY@tc##1{\textcolor[rgb]{0.24,0.48,0.48}{##1}}}
\@namedef{PY@tok@c1}{\let\PY@it=\textit\def\PY@tc##1{\textcolor[rgb]{0.24,0.48,0.48}{##1}}}
\@namedef{PY@tok@cs}{\let\PY@it=\textit\def\PY@tc##1{\textcolor[rgb]{0.24,0.48,0.48}{##1}}}

\def\PYZbs{\char`\\}
\def\PYZus{\char`\_}
\def\PYZob{\char`\{}
\def\PYZcb{\char`\}}
\def\PYZca{\char`\^}
\def\PYZam{\char`\&}
\def\PYZlt{\char`\<}
\def\PYZgt{\char`\>}
\def\PYZsh{\char`\#}
\def\PYZpc{\char`\%}
\def\PYZdl{\char`\$}
\def\PYZhy{\char`\-}
\def\PYZsq{\char`\'}
\def\PYZdq{\char`\"}
\def\PYZti{\char`\~}
% for compatibility with earlier versions
\def\PYZat{@}
\def\PYZlb{[}
\def\PYZrb{]}
\makeatother


    % For linebreaks inside Verbatim environment from package fancyvrb.
    \makeatletter
        \newbox\Wrappedcontinuationbox
        \newbox\Wrappedvisiblespacebox
        \newcommand*\Wrappedvisiblespace {\textcolor{red}{\textvisiblespace}}
        \newcommand*\Wrappedcontinuationsymbol {\textcolor{red}{\llap{\tiny$\m@th\hookrightarrow$}}}
        \newcommand*\Wrappedcontinuationindent {3ex }
        \newcommand*\Wrappedafterbreak {\kern\Wrappedcontinuationindent\copy\Wrappedcontinuationbox}
        % Take advantage of the already applied Pygments mark-up to insert
        % potential linebreaks for TeX processing.
        %        {, <, #, %, $, ' and ": go to next line.
        %        _, }, ^, &, >, - and ~: stay at end of broken line.
        % Use of \textquotesingle for straight quote.
        \newcommand*\Wrappedbreaksatspecials {%
            \def\PYGZus{\discretionary{\char`\_}{\Wrappedafterbreak}{\char`\_}}%
            \def\PYGZob{\discretionary{}{\Wrappedafterbreak\char`\{}{\char`\{}}%
            \def\PYGZcb{\discretionary{\char`\}}{\Wrappedafterbreak}{\char`\}}}%
            \def\PYGZca{\discretionary{\char`\^}{\Wrappedafterbreak}{\char`\^}}%
            \def\PYGZam{\discretionary{\char`\&}{\Wrappedafterbreak}{\char`\&}}%
            \def\PYGZlt{\discretionary{}{\Wrappedafterbreak\char`\<}{\char`\<}}%
            \def\PYGZgt{\discretionary{\char`\>}{\Wrappedafterbreak}{\char`\>}}%
            \def\PYGZsh{\discretionary{}{\Wrappedafterbreak\char`\#}{\char`\#}}%
            \def\PYGZpc{\discretionary{}{\Wrappedafterbreak\char`\%}{\char`\%}}%
            \def\PYGZdl{\discretionary{}{\Wrappedafterbreak\char`\$}{\char`\$}}%
            \def\PYGZhy{\discretionary{\char`\-}{\Wrappedafterbreak}{\char`\-}}%
            \def\PYGZsq{\discretionary{}{\Wrappedafterbreak\textquotesingle}{\textquotesingle}}%
            \def\PYGZdq{\discretionary{}{\Wrappedafterbreak\char`\"}{\char`\"}}%
            \def\PYGZti{\discretionary{\char`\~}{\Wrappedafterbreak}{\char`\~}}%
        }
        % Some characters . , ; ? ! / are not pygmentized.
        % This macro makes them "active" and they will insert potential linebreaks
        \newcommand*\Wrappedbreaksatpunct {%
            \lccode`\~`\.\lowercase{\def~}{\discretionary{\hbox{\char`\.}}{\Wrappedafterbreak}{\hbox{\char`\.}}}%
            \lccode`\~`\,\lowercase{\def~}{\discretionary{\hbox{\char`\,}}{\Wrappedafterbreak}{\hbox{\char`\,}}}%
            \lccode`\~`\;\lowercase{\def~}{\discretionary{\hbox{\char`\;}}{\Wrappedafterbreak}{\hbox{\char`\;}}}%
            \lccode`\~`\:\lowercase{\def~}{\discretionary{\hbox{\char`\:}}{\Wrappedafterbreak}{\hbox{\char`\:}}}%
            \lccode`\~`\?\lowercase{\def~}{\discretionary{\hbox{\char`\?}}{\Wrappedafterbreak}{\hbox{\char`\?}}}%
            \lccode`\~`\!\lowercase{\def~}{\discretionary{\hbox{\char`\!}}{\Wrappedafterbreak}{\hbox{\char`\!}}}%
            \lccode`\~`\/\lowercase{\def~}{\discretionary{\hbox{\char`\/}}{\Wrappedafterbreak}{\hbox{\char`\/}}}%
            \catcode`\.\active
            \catcode`\,\active
            \catcode`\;\active
            \catcode`\:\active
            \catcode`\?\active
            \catcode`\!\active
            \catcode`\/\active
            \lccode`\~`\~
        }
    \makeatother

    \let\OriginalVerbatim=\Verbatim
    \makeatletter
    \renewcommand{\Verbatim}[1][1]{%
        %\parskip\z@skip
        \sbox\Wrappedcontinuationbox {\Wrappedcontinuationsymbol}%
        \sbox\Wrappedvisiblespacebox {\FV@SetupFont\Wrappedvisiblespace}%
        \def\FancyVerbFormatLine ##1{\hsize\linewidth
            \vtop{\raggedright\hyphenpenalty\z@\exhyphenpenalty\z@
                \doublehyphendemerits\z@\finalhyphendemerits\z@
                \strut ##1\strut}%
        }%
        % If the linebreak is at a space, the latter will be displayed as visible
        % space at end of first line, and a continuation symbol starts next line.
        % Stretch/shrink are however usually zero for typewriter font.
        \def\FV@Space {%
            \nobreak\hskip\z@ plus\fontdimen3\font minus\fontdimen4\font
            \discretionary{\copy\Wrappedvisiblespacebox}{\Wrappedafterbreak}
            {\kern\fontdimen2\font}%
        }%

        % Allow breaks at special characters using \PYG... macros.
        \Wrappedbreaksatspecials
        % Breaks at punctuation characters . , ; ? ! and / need catcode=\active
        \OriginalVerbatim[#1,codes*=\Wrappedbreaksatpunct]%
    }
    \makeatother

    % Exact colors from NB
    \definecolor{incolor}{HTML}{303F9F}
    \definecolor{outcolor}{HTML}{D84315}
    \definecolor{cellborder}{HTML}{CFCFCF}
    \definecolor{cellbackground}{HTML}{F7F7F7}

    % prompt
    \makeatletter
    \newcommand{\boxspacing}{\kern\kvtcb@left@rule\kern\kvtcb@boxsep}
    \makeatother
    \newcommand{\prompt}[4]{
        {\ttfamily\llap{{\color{#2}[#3]:\hspace{3pt}#4}}\vspace{-\baselineskip}}
    }
    

    
    % Prevent overflowing lines due to hard-to-break entities
    \sloppy
    % Setup hyperref package
    \hypersetup{
      breaklinks=true,  % so long urls are correctly broken across lines
      colorlinks=true,
      urlcolor=urlcolor,
      linkcolor=linkcolor,
      citecolor=citecolor,
      }
    % Slightly bigger margins than the latex defaults
    
    \geometry{verbose,tmargin=1in,bmargin=1in,lmargin=1in,rmargin=1in}
    
    

\begin{document}
    
    \maketitle
    
    

    
    \begin{tcolorbox}[breakable, size=fbox, boxrule=1pt, pad at break*=1mm,colback=cellbackground, colframe=cellborder]
\prompt{In}{incolor}{151}{\boxspacing}
\begin{Verbatim}[commandchars=\\\{\}]
\PY{o}{\PYZpc{}}\PY{k}{matplotlib} inline

\PY{k+kn}{import} \PY{n+nn}{numpy} \PY{k}{as} \PY{n+nn}{np}
\PY{k+kn}{from} \PY{n+nn}{matplotlib} \PY{k+kn}{import} \PY{n}{pyplot} \PY{k}{as} \PY{n}{plt}
\PY{k+kn}{from} \PY{n+nn}{qiskit} \PY{k+kn}{import} \PY{n}{QuantumCircuit}
\PY{k+kn}{from} \PY{n+nn}{qiskit}\PY{n+nn}{.}\PY{n+nn}{primitives} \PY{k+kn}{import} \PY{n}{StatevectorSampler}
\PY{k+kn}{from} \PY{n+nn}{qiskit}\PY{n+nn}{.}\PY{n+nn}{quantum\PYZus{}info} \PY{k+kn}{import} \PY{n}{Operator}
\PY{k+kn}{from} \PY{n+nn}{qiskit}\PY{n+nn}{.}\PY{n+nn}{quantum\PYZus{}info} \PY{k+kn}{import} \PY{n}{Statevector}

\PY{n}{sampler} \PY{o}{=} \PY{n}{StatevectorSampler}\PY{p}{(}\PY{p}{)}
\end{Verbatim}
\end{tcolorbox}

    \hypertarget{iiqc-trabalho-2-miguel-curto-castela-uc2022212972}{%
\subsubsection{\texorpdfstring{IIQC Trabalho 2, Miguel Curto Castela,
uc2022212972
}{IIQC Trabalho 2, Miguel Curto Castela, uc2022212972 }}\label{iiqc-trabalho-2-miguel-curto-castela-uc2022212972}}

    \hypertarget{exercuxedcios-simples-1}{%
\subsubsection{\texorpdfstring{2.1 Exercícios simples 1
}{2.1 Exercícios simples 1 }}\label{exercuxedcios-simples-1}}

\begin{enumerate}
\def\labelenumi{\arabic{enumi}.}
\tightlist
\item
  Escrever o produto Kronecker (produto tensorial) dos qubits:

  \begin{enumerate}
  \def\labelenumii{\alph{enumii})}
  \tightlist
  \item
    \(|0\rangle\otimes|1\rangle\)\\
  \item
    \(|0\rangle\otimes|+\rangle\)\\
  \item
    \(|+\rangle\otimes|1\rangle\)\\
  \item
    \(|-\rangle\otimes|+\rangle\)
  \end{enumerate}
\end{enumerate}

tendo em conta que:

\[ |+\rangle = \tfrac{1}{\sqrt{2}}(|0\rangle + |1\rangle) = \tfrac{1}{\sqrt{2}}\begin{bmatrix} 1 \\ 1 \end{bmatrix}\]
\[ |-\rangle = \tfrac{1}{\sqrt{2}}(|0\rangle - |1\rangle) = \tfrac{1}{\sqrt{2}}\begin{bmatrix} 1 \\ -1 \end{bmatrix} \]

    Resposta ao exercício no PDF exercicios\_manuscritos\_2022212972.pdf
incluído na entrega

    \begin{enumerate}
\def\labelenumi{\arabic{enumi}.}
\setcounter{enumi}{1}
\tightlist
\item
  Escreva o estado: \$\textbar{}\psi\rangle =
  \tfrac{1}{\sqrt{2}}\textbar00\rangle +
  \tfrac{i}{\sqrt{2}}\textbar01\rangle \$ como dois qubits seperarados.
\end{enumerate}

    Resposta ao exercício no PDF exercicios\_manuscritos\_2022212972.pdf
incluído na entrega

    \hypertarget{exercuxedcios-simples-2}{%
\subsubsection{\texorpdfstring{2.2 Exercícios simples 2
}{2.2 Exercícios simples 2 }}\label{exercuxedcios-simples-2}}

\begin{enumerate}
\def\labelenumi{\arabic{enumi}.}
\tightlist
\item
  Calcule a unitária (\(U\)) de um circuito de um quibit criado pela
  sequência de portas: \(U = XZH\). Utilize o Qiskit para verificar os
  seus resultados.
\end{enumerate}

    Resposta ao exercício no PDF exercicios\_manuscritos\_2022212972.pdf
incluído na entrega

    \begin{tcolorbox}[breakable, size=fbox, boxrule=1pt, pad at break*=1mm,colback=cellbackground, colframe=cellborder]
\prompt{In}{incolor}{152}{\boxspacing}
\begin{Verbatim}[commandchars=\\\{\}]
\PY{c+c1}{\PYZsh{}verificação feita com Qiskit}
\PY{n}{qc} \PY{o}{=} \PY{n}{QuantumCircuit}\PY{p}{(}\PY{l+m+mi}{1}\PY{p}{)}
\PY{n}{qc}\PY{o}{.}\PY{n}{x}\PY{p}{(}\PY{l+m+mi}{0}\PY{p}{)}
\PY{n}{qc}\PY{o}{.}\PY{n}{z}\PY{p}{(}\PY{l+m+mi}{0}\PY{p}{)}
\PY{n}{qc}\PY{o}{.}\PY{n}{h}\PY{p}{(}\PY{l+m+mi}{0}\PY{p}{)}

\PY{n}{unitary} \PY{o}{=} \PY{n}{Operator}\PY{p}{(}\PY{n}{qc}\PY{p}{)}\PY{o}{.}\PY{n}{data}
\PY{n+nb}{print}\PY{p}{(}\PY{l+s+s2}{\PYZdq{}}\PY{l+s+s2}{unitary matrix:}\PY{l+s+s2}{\PYZdq{}}\PY{p}{)}
\PY{n+nb}{print}\PY{p}{(}\PY{n}{unitary}\PY{p}{)}
\end{Verbatim}
\end{tcolorbox}

    \begin{Verbatim}[commandchars=\\\{\}]
unitary matrix:
[[-0.70710678+0.j  0.70710678+0.j]
 [ 0.70710678+0.j  0.70710678+0.j]]
    \end{Verbatim}

    \begin{enumerate}
\def\labelenumi{\arabic{enumi}.}
\setcounter{enumi}{1}
\tightlist
\item
  Considere um circuito diferente com as mesmas portas, mas
  ``empilhadas'' para ter um circuito com 3 qubits, aplicando cada porta
  a um qubit diferente. Calcule o seu produto kronecker (produto
  tensorial) para obter a unitária (\(U\)) do circuito de 3 quibits, em
  seguida, verifique a sua resposta utilizando o Qiskit.
\end{enumerate}

    Resposta ao exercício no PDF exercicios\_manuscritos\_2022212972.pdf
incluído na entrega

    \begin{tcolorbox}[breakable, size=fbox, boxrule=1pt, pad at break*=1mm,colback=cellbackground, colframe=cellborder]
\prompt{In}{incolor}{153}{\boxspacing}
\begin{Verbatim}[commandchars=\\\{\}]
\PY{c+c1}{\PYZsh{}verificação feita com Qiskit}
\PY{n}{qc} \PY{o}{=} \PY{n}{QuantumCircuit}\PY{p}{(}\PY{l+m+mi}{3}\PY{p}{)}
\PY{n}{qc}\PY{o}{.}\PY{n}{x}\PY{p}{(}\PY{l+m+mi}{2}\PY{p}{)}
\PY{n}{qc}\PY{o}{.}\PY{n}{z}\PY{p}{(}\PY{l+m+mi}{1}\PY{p}{)}
\PY{n}{qc}\PY{o}{.}\PY{n}{h}\PY{p}{(}\PY{l+m+mi}{0}\PY{p}{)}

\PY{n}{unitary} \PY{o}{=} \PY{n}{Operator}\PY{p}{(}\PY{n}{qc}\PY{p}{)}\PY{o}{.}\PY{n}{data}
\PY{n+nb}{print}\PY{p}{(}\PY{l+s+s2}{\PYZdq{}}\PY{l+s+s2}{unitary matrix:}\PY{l+s+s2}{\PYZdq{}}\PY{p}{)}
\PY{n+nb}{print}\PY{p}{(}\PY{n}{unitary}\PY{p}{)}
\end{Verbatim}
\end{tcolorbox}

    \begin{Verbatim}[commandchars=\\\{\}]
unitary matrix:
[[ 0.        +0.j  0.        +0.j  0.        +0.j  0.        +0.j
   0.70710678+0.j  0.70710678+0.j  0.        +0.j  0.        +0.j]
 [ 0.        +0.j  0.        +0.j  0.        +0.j  0.        +0.j
   0.70710678+0.j -0.70710678+0.j  0.        +0.j  0.        +0.j]
 [ 0.        +0.j  0.        +0.j  0.        +0.j  0.        +0.j
   0.        +0.j  0.        +0.j -0.70710678+0.j -0.70710678+0.j]
 [ 0.        +0.j  0.        +0.j  0.        +0.j  0.        +0.j
   0.        +0.j  0.        +0.j -0.70710678+0.j  0.70710678+0.j]
 [ 0.70710678+0.j  0.70710678+0.j  0.        +0.j  0.        +0.j
   0.        +0.j  0.        +0.j  0.        +0.j  0.        +0.j]
 [ 0.70710678+0.j -0.70710678+0.j  0.        +0.j  0.        +0.j
   0.        +0.j  0.        +0.j  0.        +0.j  0.        +0.j]
 [ 0.        +0.j  0.        +0.j -0.70710678+0.j -0.70710678+0.j
   0.        +0.j  0.        +0.j  0.        +0.j  0.        +0.j]
 [ 0.        +0.j  0.        +0.j -0.70710678+0.j  0.70710678+0.j
   0.        +0.j  0.        +0.j  0.        +0.j  0.        +0.j]]
    \end{Verbatim}

    \hypertarget{exercuxedcios-simples-3}{%
\subsubsection{\texorpdfstring{Exercícios simples 3
}{Exercícios simples 3 }}\label{exercuxedcios-simples-3}}

\begin{enumerate}
\def\labelenumi{\arabic{enumi}.}
\tightlist
\item
  Crie um circuito quântico que produza o estado de Bell:
  \(\tfrac{1}{\sqrt{2}}(|01\rangle + |10\rangle)\). Utilize o Qiskit
  para construir o circuito e statevector para verificar o seu
  resultado.
\end{enumerate}

    \begin{tcolorbox}[breakable, size=fbox, boxrule=1pt, pad at break*=1mm,colback=cellbackground, colframe=cellborder]
\prompt{In}{incolor}{154}{\boxspacing}
\begin{Verbatim}[commandchars=\\\{\}]
\PY{c+c1}{\PYZsh{}Bell state: |psi⟩ = 1/sqrt(2) (|01⟩ + |10⟩)}
\PY{c+c1}{\PYZsh{}Circuit to create the Bell state}
\PY{n}{bell\PYZus{}state\PYZus{}expected} \PY{o}{=} \PY{n}{Statevector}\PY{p}{(}\PY{p}{[}\PY{l+m+mi}{0}\PY{p}{,} \PY{l+m+mi}{1}\PY{p}{,} \PY{l+m+mi}{1}\PY{p}{,} \PY{l+m+mi}{0}\PY{p}{]}\PY{p}{)} \PY{o}{/} \PY{n}{np}\PY{o}{.}\PY{n}{sqrt}\PY{p}{(}\PY{l+m+mi}{2}\PY{p}{)}

\PY{n}{qc} \PY{o}{=} \PY{n}{QuantumCircuit}\PY{p}{(}\PY{l+m+mi}{2}\PY{p}{)}
\PY{n}{qc}\PY{o}{.}\PY{n}{x}\PY{p}{(}\PY{l+m+mi}{1}\PY{p}{)}
\PY{n}{qc}\PY{o}{.}\PY{n}{h}\PY{p}{(}\PY{l+m+mi}{0}\PY{p}{)}
\PY{n}{qc}\PY{o}{.}\PY{n}{cx}\PY{p}{(}\PY{l+m+mi}{0}\PY{p}{,} \PY{l+m+mi}{1}\PY{p}{)}

\PY{n+nb}{print}\PY{p}{(}\PY{n}{qc}\PY{p}{)}

\PY{n}{statevector} \PY{o}{=} \PY{n}{Statevector}\PY{o}{.}\PY{n}{from\PYZus{}instruction}\PY{p}{(}\PY{n}{qc}\PY{p}{)}
\PY{n+nb}{print}\PY{p}{(}\PY{l+s+s2}{\PYZdq{}}\PY{l+s+s2}{circuit statevector}\PY{l+s+s2}{\PYZdq{}}\PY{p}{,} \PY{n}{statevector}\PY{p}{)}

\PY{k}{if} \PY{n}{statevector}\PY{o}{.}\PY{n}{equiv}\PY{p}{(}\PY{n}{bell\PYZus{}state\PYZus{}expected}\PY{p}{)}\PY{p}{:}
    \PY{n+nb}{print}\PY{p}{(}\PY{l+s+s2}{\PYZdq{}}\PY{l+s+s2}{O circuito produz o bell state esperado}\PY{l+s+s2}{\PYZdq{}}\PY{p}{)}
\PY{k}{else}\PY{p}{:}
    \PY{n+nb}{print}\PY{p}{(}\PY{l+s+s2}{\PYZdq{}}\PY{l+s+s2}{O circuito não produz o bell state esperado}\PY{l+s+s2}{\PYZdq{}}\PY{p}{)}
\end{Verbatim}
\end{tcolorbox}

    \begin{Verbatim}[commandchars=\\\{\}]
     ┌───┐
q\_0: ┤ H ├──■──
     ├───┤┌─┴─┐
q\_1: ┤ X ├┤ X ├
     └───┘└───┘
circuit statevector Statevector([0.        +0.j, 0.70710678+0.j, 0.70710678+0.j,
             0.        +0.j],
            dims=(2, 2))
O circuito produz o bell state esperado
    \end{Verbatim}

    \begin{enumerate}
\def\labelenumi{\arabic{enumi}.}
\setcounter{enumi}{1}
\tightlist
\item
  O circuito que criou acima transforma o estado \(|00\rangle\) em
  \(\tfrac{1}{\sqrt{2}}(|01\rangle + |10\rangle)\), calcule a matriz
  unitária deste circuito utilizando o Qiskit. Verifique se esta
  unitária realiza de facto a transformação correta.
\end{enumerate}

    \begin{tcolorbox}[breakable, size=fbox, boxrule=1pt, pad at break*=1mm,colback=cellbackground, colframe=cellborder]
\prompt{In}{incolor}{155}{\boxspacing}
\begin{Verbatim}[commandchars=\\\{\}]
\PY{c+c1}{\PYZsh{}Matriz unitária do circuito}
\PY{n}{unit} \PY{o}{=} \PY{n}{Operator}\PY{p}{(}\PY{n}{qc}\PY{p}{)}\PY{o}{.}\PY{n}{data}
\PY{n+nb}{print}\PY{p}{(}\PY{l+s+s2}{\PYZdq{}}\PY{l+s+s2}{unitary matrix:}\PY{l+s+s2}{\PYZdq{}}\PY{p}{)}
\PY{n+nb}{print}\PY{p}{(}\PY{n}{unit}\PY{p}{)}

\PY{n}{initial\PYZus{}state} \PY{o}{=} \PY{n}{Statevector}\PY{o}{.}\PY{n}{from\PYZus{}label}\PY{p}{(}\PY{l+s+s1}{\PYZsq{}}\PY{l+s+s1}{00}\PY{l+s+s1}{\PYZsq{}}\PY{p}{)}


\PY{n}{final\PYZus{}state} \PY{o}{=} \PY{n}{initial\PYZus{}state}\PY{o}{.}\PY{n}{evolve}\PY{p}{(}\PY{n}{unit}\PY{p}{)}
\PY{n+nb}{print}\PY{p}{(}\PY{l+s+s2}{\PYZdq{}}\PY{l+s+s2}{final state:}\PY{l+s+s2}{\PYZdq{}}\PY{p}{,}\PY{n}{final\PYZus{}state}\PY{p}{)}

\PY{k}{if} \PY{n}{statevector}\PY{o}{.}\PY{n}{equiv}\PY{p}{(}\PY{n}{bell\PYZus{}state\PYZus{}expected}\PY{p}{)}\PY{p}{:}
    \PY{n+nb}{print}\PY{p}{(}\PY{l+s+s2}{\PYZdq{}}\PY{l+s+s2}{A matriz unitária produz o bell state esperado}\PY{l+s+s2}{\PYZdq{}}\PY{p}{)}
\PY{k}{else}\PY{p}{:}
    \PY{n+nb}{print}\PY{p}{(}\PY{l+s+s2}{\PYZdq{}}\PY{l+s+s2}{A matriz unitária não produz o bell state esperado}\PY{l+s+s2}{\PYZdq{}}\PY{p}{)}
\end{Verbatim}
\end{tcolorbox}

    \begin{Verbatim}[commandchars=\\\{\}]
unitary matrix:
[[ 0.        +0.j  0.        +0.j  0.70710678+0.j  0.70710678+0.j]
 [ 0.70710678+0.j -0.70710678+0.j  0.        +0.j  0.        +0.j]
 [ 0.70710678+0.j  0.70710678+0.j  0.        +0.j  0.        +0.j]
 [ 0.        +0.j  0.        +0.j  0.70710678+0.j -0.70710678+0.j]]
final state: Statevector([0.        +0.j, 0.70710678+0.j, 0.70710678+0.j,
             0.        +0.j],
            dims=(2, 2))
A matriz unitária produz o bell state esperado
    \end{Verbatim}

    \hypertarget{algoritmo-deutsch-jozsa}{%
\subsubsection{\texorpdfstring{2.4 Algoritmo Deutsch-Jozsa
}{2.4 Algoritmo Deutsch-Jozsa }}\label{algoritmo-deutsch-jozsa}}

Apresente a sua implementação, simulação de alguns casos, e execução em
hardware real de pelo menos dois casos aproventando o número de qubits
disponiveis.

    \begin{tcolorbox}[breakable, size=fbox, boxrule=1pt, pad at break*=1mm,colback=cellbackground, colframe=cellborder]
\prompt{In}{incolor}{156}{\boxspacing}
\begin{Verbatim}[commandchars=\\\{\}]
\PY{k+kn}{import} \PY{n+nn}{numpy} \PY{k}{as} \PY{n+nn}{np}
\PY{k+kn}{import} \PY{n+nn}{matplotlib}\PY{n+nn}{.}\PY{n+nn}{pyplot} \PY{k}{as} \PY{n+nn}{plt}
\PY{k+kn}{from} \PY{n+nn}{qiskit}\PY{n+nn}{.}\PY{n+nn}{visualization} \PY{k+kn}{import} \PY{n}{plot\PYZus{}histogram}
\PY{k+kn}{from} \PY{n+nn}{qiskit} \PY{k+kn}{import} \PY{n}{QuantumCircuit}\PY{p}{,} \PY{n}{transpile}
\PY{k+kn}{from} \PY{n+nn}{qiskit\PYZus{}aer} \PY{k+kn}{import} \PY{n}{Aer}
\PY{k+kn}{from} \PY{n+nn}{qiskit}\PY{n+nn}{.}\PY{n+nn}{visualization} \PY{k+kn}{import} \PY{n}{plot\PYZus{}histogram}
\PY{k+kn}{from} \PY{n+nn}{qiskit\PYZus{}aer} \PY{k+kn}{import} \PY{n}{AerSimulator}
\PY{k+kn}{from} \PY{n+nn}{qiskit\PYZus{}ibm\PYZus{}runtime} \PY{k+kn}{import} \PY{n}{QiskitRuntimeService}
\PY{k+kn}{import} \PY{n+nn}{os}
\PY{k+kn}{from} \PY{n+nn}{dotenv} \PY{k+kn}{import} \PY{n}{load\PYZus{}dotenv}

\PY{k}{if} \PY{l+s+s2}{\PYZdq{}}\PY{l+s+s2}{IBM\PYZus{}TOKEN}\PY{l+s+s2}{\PYZdq{}} \PY{o+ow}{in} \PY{n}{os}\PY{o}{.}\PY{n}{environ}\PY{p}{:}
    \PY{k}{del} \PY{n}{os}\PY{o}{.}\PY{n}{environ}\PY{p}{[}\PY{l+s+s2}{\PYZdq{}}\PY{l+s+s2}{IBM\PYZus{}TOKEN}\PY{l+s+s2}{\PYZdq{}}\PY{p}{]}
\PY{n}{load\PYZus{}dotenv}\PY{p}{(}\PY{p}{)}
\PY{n}{IBM\PYZus{}TOKEN} \PY{o}{=} \PY{n}{os}\PY{o}{.}\PY{n}{getenv}\PY{p}{(}\PY{l+s+s2}{\PYZdq{}}\PY{l+s+s2}{IBM\PYZus{}TOKEN}\PY{l+s+s2}{\PYZdq{}}\PY{p}{)}
\PY{c+c1}{\PYZsh{}print(f\PYZdq{}IBM\PYZus{}TOKEN: \PYZob{}IBM\PYZus{}TOKEN\PYZcb{}\PYZdq{})}
\PY{n}{service} \PY{o}{=} \PY{n}{QiskitRuntimeService}\PY{p}{(}\PY{n}{channel}\PY{o}{=}\PY{l+s+s2}{\PYZdq{}}\PY{l+s+s2}{ibm\PYZus{}quantum}\PY{l+s+s2}{\PYZdq{}}\PY{p}{,} \PY{n}{token}\PY{o}{=}\PY{n}{IBM\PYZus{}TOKEN}\PY{p}{)}
\PY{n}{backend} \PY{o}{=} \PY{n}{service}\PY{o}{.}\PY{n}{least\PYZus{}busy}\PY{p}{(}\PY{n}{operational}\PY{o}{=}\PY{k+kc}{True}\PY{p}{,} \PY{n}{simulator}\PY{o}{=}\PY{k+kc}{False}\PY{p}{)}
\PY{n}{sim} \PY{o}{=} \PY{n}{AerSimulator}\PY{p}{(}\PY{p}{)}
\PY{n+nb}{print}\PY{p}{(}\PY{n}{backend}\PY{p}{)}
\end{Verbatim}
\end{tcolorbox}

    \begin{Verbatim}[commandchars=\\\{\}]
<IBMBackend('ibm\_kyiv')>
    \end{Verbatim}

    \begin{tcolorbox}[breakable, size=fbox, boxrule=1pt, pad at break*=1mm,colback=cellbackground, colframe=cellborder]
\prompt{In}{incolor}{157}{\boxspacing}
\begin{Verbatim}[commandchars=\\\{\}]
\PY{k}{def} \PY{n+nf}{constantOracle}\PY{p}{(}\PY{n+nb}{len}\PY{p}{)}\PY{p}{:}
    \PY{n}{oracle} \PY{o}{=} \PY{n}{QuantumCircuit}\PY{p}{(}\PY{n+nb}{len}\PY{o}{+}\PY{l+m+mi}{1}\PY{p}{)}
    
    \PY{n}{oracle}\PY{o}{.}\PY{n}{barrier}\PY{p}{(}\PY{p}{)}
    
    \PY{n}{output} \PY{o}{=} \PY{n}{np}\PY{o}{.}\PY{n}{random}\PY{o}{.}\PY{n}{randint}\PY{p}{(}\PY{l+m+mi}{2}\PY{p}{)}
    \PY{k}{if}\PY{p}{(}\PY{n}{output} \PY{o}{==} \PY{l+m+mi}{1}\PY{p}{)}\PY{p}{:}
        \PY{n}{oracle}\PY{o}{.}\PY{n}{x}\PY{p}{(}\PY{n+nb}{len}\PY{p}{)}

    \PY{n}{oracle}\PY{o}{.}\PY{n}{barrier}\PY{p}{(}\PY{p}{)}

    \PY{k}{return} \PY{n}{oracle}

\PY{k}{def} \PY{n+nf}{balancedOracle}\PY{p}{(}\PY{n+nb}{len}\PY{p}{)}\PY{p}{:}
    \PY{n}{oracle} \PY{o}{=} \PY{n}{QuantumCircuit}\PY{p}{(}\PY{n+nb}{len}\PY{o}{+}\PY{l+m+mi}{1}\PY{p}{)}
    
    \PY{n}{oracle}\PY{o}{.}\PY{n}{barrier}\PY{p}{(}\PY{p}{)}

    \PY{c+c1}{\PYZsh{}Random implementation}
    \PY{n}{b} \PY{o}{=} \PY{n}{np}\PY{o}{.}\PY{n}{random}\PY{o}{.}\PY{n}{randint}\PY{p}{(}\PY{l+m+mi}{1}\PY{p}{,} \PY{l+m+mi}{2}\PY{o}{*}\PY{o}{*}\PY{n+nb}{len}\PY{p}{)}
    \PY{n}{b\PYZus{}str} \PY{o}{=} \PY{n+nb}{format}\PY{p}{(}\PY{n}{b}\PY{p}{,} \PY{l+s+s1}{\PYZsq{}}\PY{l+s+s1}{0}\PY{l+s+s1}{\PYZsq{}}\PY{o}{+}\PY{n+nb}{str}\PY{p}{(}\PY{n+nb}{len}\PY{p}{)}\PY{o}{+}\PY{l+s+s1}{\PYZsq{}}\PY{l+s+s1}{b}\PY{l+s+s1}{\PYZsq{}}\PY{p}{)}
    
    \PY{k}{for} \PY{n}{qb} \PY{o+ow}{in} \PY{n+nb}{range}\PY{p}{(}\PY{n+nb}{len}\PY{p}{)}\PY{p}{:}
        \PY{k}{if}\PY{p}{(}\PY{n}{b\PYZus{}str}\PY{p}{[}\PY{n}{qb}\PY{p}{]} \PY{o}{==} \PY{l+s+s1}{\PYZsq{}}\PY{l+s+s1}{1}\PY{l+s+s1}{\PYZsq{}}\PY{p}{)}\PY{p}{:}
            \PY{n}{oracle}\PY{o}{.}\PY{n}{x}\PY{p}{(}\PY{n}{qb}\PY{p}{)}

        
    \PY{k}{for} \PY{n}{qb} \PY{o+ow}{in} \PY{n+nb}{range}\PY{p}{(}\PY{n+nb}{len}\PY{p}{)}\PY{p}{:}
        \PY{n}{oracle}\PY{o}{.}\PY{n}{cx}\PY{p}{(}\PY{n}{qb}\PY{p}{,} \PY{n+nb}{len}\PY{p}{)}
        
    \PY{k}{for} \PY{n}{qb} \PY{o+ow}{in} \PY{n+nb}{range}\PY{p}{(}\PY{n+nb}{len}\PY{p}{)}\PY{p}{:}
        \PY{k}{if}\PY{p}{(}\PY{n}{b\PYZus{}str}\PY{p}{[}\PY{n}{qb}\PY{p}{]} \PY{o}{==} \PY{l+s+s1}{\PYZsq{}}\PY{l+s+s1}{1}\PY{l+s+s1}{\PYZsq{}}\PY{p}{)}\PY{p}{:}
            \PY{n}{oracle}\PY{o}{.}\PY{n}{x}\PY{p}{(}\PY{n}{qb}\PY{p}{)}    

    \PY{n}{oracle}\PY{o}{.}\PY{n}{barrier}\PY{p}{(}\PY{p}{)}

    \PY{c+c1}{\PYZsh{}Static implementation }
\PY{+w}{    }\PY{l+s+sd}{\PYZsq{}\PYZsq{}\PYZsq{}}
\PY{l+s+sd}{    oracle = QuantumCircuit(len+1)}

\PY{l+s+sd}{    oracle.barrier()}

\PY{l+s+sd}{    for qb in range(len):}
\PY{l+s+sd}{        if(qb \PYZpc{} 2 == 0):}
\PY{l+s+sd}{            oracle.x(qb)}

\PY{l+s+sd}{    oracle.barrier()}

\PY{l+s+sd}{    for qb in range(len):}
\PY{l+s+sd}{        oracle.cx(qb, len)}

\PY{l+s+sd}{    oracle.barrier()}

\PY{l+s+sd}{    for qb in range(len):}
\PY{l+s+sd}{        if(qb \PYZpc{} 2 == 0):}
\PY{l+s+sd}{            oracle.x(qb)    }

\PY{l+s+sd}{    oracle.barrier()}
\PY{l+s+sd}{    \PYZsq{}\PYZsq{}\PYZsq{}}

    \PY{k}{return} \PY{n}{oracle}
\end{Verbatim}
\end{tcolorbox}

    \begin{tcolorbox}[breakable, size=fbox, boxrule=1pt, pad at break*=1mm,colback=cellbackground, colframe=cellborder]
\prompt{In}{incolor}{158}{\boxspacing}
\begin{Verbatim}[commandchars=\\\{\}]
\PY{k}{def} \PY{n+nf}{createOracle}\PY{p}{(}\PY{n+nb}{len}\PY{p}{)}\PY{p}{:}
    \PY{n}{constant} \PY{o}{=} \PY{n}{constantOracle}\PY{p}{(}\PY{n+nb}{len}\PY{p}{)}
    \PY{n}{balanced} \PY{o}{=} \PY{n}{balancedOracle}\PY{p}{(}\PY{n+nb}{len}\PY{p}{)}
    
    \PY{k}{return} \PY{n}{constant}\PY{p}{,} \PY{n}{balanced}
\end{Verbatim}
\end{tcolorbox}

    \begin{tcolorbox}[breakable, size=fbox, boxrule=1pt, pad at break*=1mm,colback=cellbackground, colframe=cellborder]
\prompt{In}{incolor}{159}{\boxspacing}
\begin{Verbatim}[commandchars=\\\{\}]
\PY{k}{def} \PY{n+nf}{createCircuit\PYZus{}dj}\PY{p}{(}\PY{n+nb}{len}\PY{p}{,} \PY{n}{oracle}\PY{p}{)}\PY{p}{:}
    \PY{n}{dj} \PY{o}{=} \PY{n}{QuantumCircuit}\PY{p}{(}\PY{n+nb}{len}\PY{o}{+}\PY{l+m+mi}{1}\PY{p}{,} \PY{n+nb}{len}\PY{p}{)}
    
    \PY{c+c1}{\PYZsh{}input qubits}
    \PY{k}{for} \PY{n}{qb} \PY{o+ow}{in} \PY{n+nb}{range}\PY{p}{(}\PY{n+nb}{len}\PY{p}{)}\PY{p}{:}
        \PY{n}{dj}\PY{o}{.}\PY{n}{h}\PY{p}{(}\PY{n}{qb}\PY{p}{)}
        
    \PY{c+c1}{\PYZsh{}output qubit}
    \PY{n}{dj}\PY{o}{.}\PY{n}{x}\PY{p}{(}\PY{n+nb}{len}\PY{p}{)}
    \PY{n}{dj}\PY{o}{.}\PY{n}{h}\PY{p}{(}\PY{n+nb}{len}\PY{p}{)}
    
    \PY{c+c1}{\PYZsh{}apply the oracle}
    \PY{n}{dj} \PY{o}{=} \PY{n}{dj}\PY{o}{.}\PY{n}{compose}\PY{p}{(}\PY{n}{oracle}\PY{p}{)}
    
    \PY{c+c1}{\PYZsh{}apply Hadamard gates to the first n qubits}
    \PY{k}{for} \PY{n}{qb} \PY{o+ow}{in} \PY{n+nb}{range}\PY{p}{(}\PY{n+nb}{len}\PY{p}{)}\PY{p}{:}
        \PY{n}{dj}\PY{o}{.}\PY{n}{h}\PY{p}{(}\PY{n}{qb}\PY{p}{)}
        
    \PY{c+c1}{\PYZsh{}measure the first n qubits}
    \PY{k}{for} \PY{n}{qb} \PY{o+ow}{in} \PY{n+nb}{range}\PY{p}{(}\PY{n+nb}{len}\PY{p}{)}\PY{p}{:}
        \PY{n}{dj}\PY{o}{.}\PY{n}{measure}\PY{p}{(}\PY{n}{qb}\PY{p}{,} \PY{n}{qb}\PY{p}{)}
    
    \PY{k}{return} \PY{n}{dj}    

\PY{k}{def} \PY{n+nf}{simulateCircuit\PYZus{}dj}\PY{p}{(}\PY{n}{dj}\PY{p}{)}\PY{p}{:}
    \PY{n}{job} \PY{o}{=} \PY{n}{sim}\PY{o}{.}\PY{n}{run}\PY{p}{(}\PY{n}{dj}\PY{p}{)}
    \PY{n}{results} \PY{o}{=} \PY{n}{job}\PY{o}{.}\PY{n}{result}\PY{p}{(}\PY{p}{)}
    \PY{n}{graph} \PY{o}{=} \PY{n}{results}\PY{o}{.}\PY{n}{get\PYZus{}counts}\PY{p}{(}\PY{p}{)}
    \PY{n}{plot\PYZus{}histogram}\PY{p}{(}\PY{n}{graph}\PY{p}{)}

\PY{k}{def} \PY{n+nf}{runCircuit\PYZus{}dj}\PY{p}{(}\PY{n}{dj}\PY{p}{)}\PY{p}{:}
    \PY{n}{transpiledCircuit} \PY{o}{=} \PY{n}{transpile}\PY{p}{(}\PY{n}{dj}\PY{p}{,} \PY{n}{backend}\PY{p}{,} \PY{n}{optimization\PYZus{}level}\PY{o}{=}\PY{l+m+mi}{3}\PY{p}{)}
    \PY{n}{transpiledCircuit}\PY{o}{.}\PY{n}{draw}\PY{p}{(}\PY{l+s+s1}{\PYZsq{}}\PY{l+s+s1}{mpl}\PY{l+s+s1}{\PYZsq{}}\PY{p}{,} \PY{n}{idle\PYZus{}wires}\PY{o}{=}\PY{k+kc}{False}\PY{p}{)}
    
    \PY{n}{job} \PY{o}{=} \PY{n}{backend}\PY{o}{.}\PY{n}{run}\PY{p}{(}\PY{p}{[}\PY{n}{transpiledCircuit}\PY{p}{]}\PY{p}{,} \PY{n}{shots}\PY{o}{=}\PY{l+m+mi}{1024}\PY{p}{)}
    \PY{n}{job}\PY{o}{.}\PY{n}{job\PYZus{}id}\PY{p}{(}\PY{p}{)}
    
    \PY{n}{results} \PY{o}{=} \PY{n}{job}\PY{o}{.}\PY{n}{result}\PY{p}{(}\PY{p}{)}
    \PY{n}{graph} \PY{o}{=} \PY{n}{results}\PY{o}{.}\PY{n}{get\PYZus{}counts}\PY{p}{(}\PY{p}{)}
    
    \PY{n}{plot\PYZus{}histogram}\PY{p}{(}\PY{n}{graph}\PY{p}{)}
    
\PY{k}{def} \PY{n+nf}{fetch\PYZus{}job\PYZus{}results}\PY{p}{(}\PY{n}{job\PYZus{}id}\PY{p}{)}\PY{p}{:}
    \PY{n}{job} \PY{o}{=} \PY{n}{service}\PY{o}{.}\PY{n}{job}\PY{p}{(}\PY{n}{job\PYZus{}id}\PY{p}{)}
    \PY{n}{results} \PY{o}{=} \PY{n}{job}\PY{o}{.}\PY{n}{result}\PY{p}{(}\PY{p}{)}
    \PY{n}{answer} \PY{o}{=} \PY{n}{results}\PY{o}{.}\PY{n}{get\PYZus{}counts}\PY{p}{(}\PY{p}{)}
    \PY{n}{plot\PYZus{}histogram}\PY{p}{(}\PY{n}{answer}\PY{p}{)}
    \PY{n}{plt}\PY{o}{.}\PY{n}{show}\PY{p}{(}\PY{p}{)}

\PY{c+c1}{\PYZsh{}número de qubits (no Deutsh\PYZhy{}Josza e no grover)}
\PY{n+nb}{len} \PY{o}{=} \PY{l+m+mi}{4}

\PY{c+c1}{\PYZsh{}create the oracles}
\PY{n}{constantOracle}\PY{p}{,} \PY{n}{balancedOracle} \PY{o}{=} \PY{n}{createOracle}\PY{p}{(}\PY{n+nb}{len}\PY{p}{)}
\end{Verbatim}
\end{tcolorbox}

    \begin{tcolorbox}[breakable, size=fbox, boxrule=1pt, pad at break*=1mm,colback=cellbackground, colframe=cellborder]
\prompt{In}{incolor}{160}{\boxspacing}
\begin{Verbatim}[commandchars=\\\{\}]
\PY{c+c1}{\PYZsh{}Deutsch\PYZhy{}Jozsa circuit with balanced oracle}
\PY{n}{balancedDj} \PY{o}{=} \PY{n}{createCircuit\PYZus{}dj}\PY{p}{(}\PY{n+nb}{len}\PY{p}{,} \PY{n}{balancedOracle}\PY{p}{)}

\PY{c+c1}{\PYZsh{}run it on real hardware}
\PY{c+c1}{\PYZsh{}runCircuit\PYZus{}dj(balancedDj)}

\PY{n+nb}{print}\PY{p}{(}\PY{l+s+s2}{\PYZdq{}}\PY{l+s+s2}{Deutch\PYZhy{}Jozsa circuit with balanced oracle }\PY{l+s+s2}{\PYZdq{}}\PY{p}{)}
\PY{n}{balancedDj}\PY{o}{.}\PY{n}{draw}\PY{p}{(}\PY{l+s+s1}{\PYZsq{}}\PY{l+s+s1}{mpl}\PY{l+s+s1}{\PYZsq{}}\PY{p}{,} \PY{n}{idle\PYZus{}wires}\PY{o}{=}\PY{k+kc}{False}\PY{p}{)}
\end{Verbatim}
\end{tcolorbox}

    \begin{Verbatim}[commandchars=\\\{\}]
Deutch-Jozsa circuit with balanced oracle
    \end{Verbatim}
 
            
\prompt{Out}{outcolor}{160}{}
    
    \begin{center}
    \adjustimage{max size={0.9\linewidth}{0.9\paperheight}}{trabalho2_files/trabalho2_21_1.png}
    \end{center}
    { \hspace*{\fill} \\}
    

    \begin{center}
    \adjustimage{max size={0.9\linewidth}{0.9\paperheight}}{trabalho2_files/trabalho2_21_2.png}
    \end{center}
    { \hspace*{\fill} \\}
    
    \begin{tcolorbox}[breakable, size=fbox, boxrule=1pt, pad at break*=1mm,colback=cellbackground, colframe=cellborder]
\prompt{In}{incolor}{161}{\boxspacing}
\begin{Verbatim}[commandchars=\\\{\}]
\PY{n+nb}{print}\PY{p}{(}\PY{l+s+s2}{\PYZdq{}}\PY{l+s+s2}{simulation of the Deutch\PYZhy{}Jozsa circuit with balanced oracle }\PY{l+s+s2}{\PYZdq{}}\PY{p}{)}
\PY{n+nb}{print}\PY{p}{(}\PY{n}{balancedDj}\PY{p}{)}
\end{Verbatim}
\end{tcolorbox}

    \begin{Verbatim}[commandchars=\\\{\}]
simulation of the Deutch-Jozsa circuit with balanced oracle
     ┌───┐      ░ ┌───┐     ┌───┐                ░ ┌───┐┌─┐
q\_0: ┤ H ├──────░─┤ X ├──■──┤ X ├────────────────░─┤ H ├┤M├─────────
     ├───┤      ░ ├───┤  │  └───┘┌───┐           ░ ├───┤└╥┘┌─┐
q\_1: ┤ H ├──────░─┤ X ├──┼────■──┤ X ├───────────░─┤ H ├─╫─┤M├──────
     ├───┤      ░ └───┘  │    │  └───┘           ░ ├───┤ ║ └╥┘┌─┐
q\_2: ┤ H ├──────░────────┼────┼────■─────────────░─┤ H ├─╫──╫─┤M├───
     ├───┤      ░ ┌───┐  │    │    │       ┌───┐ ░ ├───┤ ║  ║ └╥┘┌─┐
q\_3: ┤ H ├──────░─┤ X ├──┼────┼────┼────■──┤ X ├─░─┤ H ├─╫──╫──╫─┤M├
     ├───┤┌───┐ ░ └───┘┌─┴─┐┌─┴─┐┌─┴─┐┌─┴─┐└───┘ ░ └───┘ ║  ║  ║ └╥┘
q\_4: ┤ X ├┤ H ├─░──────┤ X ├┤ X ├┤ X ├┤ X ├──────░───────╫──╫──╫──╫─
     └───┘└───┘ ░      └───┘└───┘└───┘└───┘      ░       ║  ║  ║  ║
c: 4/════════════════════════════════════════════════════╩══╩══╩══╩═
                                                         0  1  2  3
    \end{Verbatim}

    \begin{tcolorbox}[breakable, size=fbox, boxrule=1pt, pad at break*=1mm,colback=cellbackground, colframe=cellborder]
\prompt{In}{incolor}{162}{\boxspacing}
\begin{Verbatim}[commandchars=\\\{\}]
\PY{n+nb}{print}\PY{p}{(}\PY{l+s+s2}{\PYZdq{}}\PY{l+s+s2}{results of the simulation}\PY{l+s+s2}{\PYZdq{}}\PY{p}{)}
\PY{n}{simulateCircuit\PYZus{}dj}\PY{p}{(}\PY{n}{balancedDj}\PY{p}{)}
\end{Verbatim}
\end{tcolorbox}

    \begin{Verbatim}[commandchars=\\\{\}]
results of the simulation
    \end{Verbatim}

    \begin{center}
    \adjustimage{max size={0.9\linewidth}{0.9\paperheight}}{trabalho2_files/trabalho2_23_1.png}
    \end{center}
    { \hspace*{\fill} \\}
    
    \begin{tcolorbox}[breakable, size=fbox, boxrule=1pt, pad at break*=1mm,colback=cellbackground, colframe=cellborder]
\prompt{In}{incolor}{163}{\boxspacing}
\begin{Verbatim}[commandchars=\\\{\}]
\PY{n+nb}{print}\PY{p}{(}\PY{l+s+s2}{\PYZdq{}}\PY{l+s+s2}{balanced iracle job results:}\PY{l+s+s2}{\PYZdq{}}\PY{p}{)}
\PY{n}{fetch\PYZus{}job\PYZus{}results}\PY{p}{(}\PY{l+s+s2}{\PYZdq{}}\PY{l+s+s2}{cxcrjgaztp300083372g}\PY{l+s+s2}{\PYZdq{}}\PY{p}{)}
\end{Verbatim}
\end{tcolorbox}

    \begin{Verbatim}[commandchars=\\\{\}]
balanced iracle job results:
    \end{Verbatim}

    \begin{center}
    \adjustimage{max size={0.9\linewidth}{0.9\paperheight}}{trabalho2_files/trabalho2_24_1.png}
    \end{center}
    { \hspace*{\fill} \\}
    
    \begin{tcolorbox}[breakable, size=fbox, boxrule=1pt, pad at break*=1mm,colback=cellbackground, colframe=cellborder]
\prompt{In}{incolor}{164}{\boxspacing}
\begin{Verbatim}[commandchars=\\\{\}]
\PY{c+c1}{\PYZsh{}Deutsch\PYZhy{}Jozsa circuit with the constant oracle}
\PY{n}{constantDj} \PY{o}{=} \PY{n}{createCircuit\PYZus{}dj}\PY{p}{(}\PY{n+nb}{len}\PY{p}{,} \PY{n}{constantOracle}\PY{p}{)}
\PY{c+c1}{\PYZsh{} run it on real hardware}
\PY{c+c1}{\PYZsh{}runCircuit\PYZus{}dj(constantOracle)}
\PY{n+nb}{print}\PY{p}{(}\PY{l+s+s2}{\PYZdq{}}\PY{l+s+s2}{Deutch\PYZhy{}Jozsa circuit with constant oracle }\PY{l+s+s2}{\PYZdq{}}\PY{p}{)}
\PY{n}{constantDj}\PY{o}{.}\PY{n}{draw}\PY{p}{(}\PY{l+s+s1}{\PYZsq{}}\PY{l+s+s1}{mpl}\PY{l+s+s1}{\PYZsq{}}\PY{p}{,} \PY{n}{idle\PYZus{}wires}\PY{o}{=}\PY{k+kc}{False}\PY{p}{)}
\end{Verbatim}
\end{tcolorbox}

    \begin{Verbatim}[commandchars=\\\{\}]
Deutch-Jozsa circuit with constant oracle
    \end{Verbatim}
 
            
\prompt{Out}{outcolor}{164}{}
    
    \begin{center}
    \adjustimage{max size={0.9\linewidth}{0.9\paperheight}}{trabalho2_files/trabalho2_25_1.png}
    \end{center}
    { \hspace*{\fill} \\}
    

    \begin{center}
    \adjustimage{max size={0.9\linewidth}{0.9\paperheight}}{trabalho2_files/trabalho2_25_2.png}
    \end{center}
    { \hspace*{\fill} \\}
    
    \begin{tcolorbox}[breakable, size=fbox, boxrule=1pt, pad at break*=1mm,colback=cellbackground, colframe=cellborder]
\prompt{In}{incolor}{165}{\boxspacing}
\begin{Verbatim}[commandchars=\\\{\}]
\PY{n+nb}{print}\PY{p}{(}\PY{l+s+s2}{\PYZdq{}}\PY{l+s+s2}{simulation of the Deutch\PYZhy{}Jozsa circuit with constant oracle }\PY{l+s+s2}{\PYZdq{}}\PY{p}{)}
\PY{n+nb}{print}\PY{p}{(}\PY{n}{constantDj}\PY{p}{)}
\end{Verbatim}
\end{tcolorbox}

    \begin{Verbatim}[commandchars=\\\{\}]
simulation of the Deutch-Jozsa circuit with constant oracle
     ┌───┐      ░       ░ ┌───┐┌─┐
q\_0: ┤ H ├──────░───────░─┤ H ├┤M├─────────
     ├───┤      ░       ░ ├───┤└╥┘┌─┐
q\_1: ┤ H ├──────░───────░─┤ H ├─╫─┤M├──────
     ├───┤      ░       ░ ├───┤ ║ └╥┘┌─┐
q\_2: ┤ H ├──────░───────░─┤ H ├─╫──╫─┤M├───
     ├───┤      ░       ░ ├───┤ ║  ║ └╥┘┌─┐
q\_3: ┤ H ├──────░───────░─┤ H ├─╫──╫──╫─┤M├
     ├───┤┌───┐ ░ ┌───┐ ░ └───┘ ║  ║  ║ └╥┘
q\_4: ┤ X ├┤ H ├─░─┤ X ├─░───────╫──╫──╫──╫─
     └───┘└───┘ ░ └───┘ ░       ║  ║  ║  ║
c: 4/═══════════════════════════╩══╩══╩══╩═
                                0  1  2  3
    \end{Verbatim}

    \begin{tcolorbox}[breakable, size=fbox, boxrule=1pt, pad at break*=1mm,colback=cellbackground, colframe=cellborder]
\prompt{In}{incolor}{166}{\boxspacing}
\begin{Verbatim}[commandchars=\\\{\}]
\PY{n+nb}{print}\PY{p}{(}\PY{l+s+s2}{\PYZdq{}}\PY{l+s+s2}{results of the simulation}\PY{l+s+s2}{\PYZdq{}}\PY{p}{)}
\PY{n}{simulateCircuit\PYZus{}dj}\PY{p}{(}\PY{n}{constantDj}\PY{p}{)}
\end{Verbatim}
\end{tcolorbox}

    \begin{Verbatim}[commandchars=\\\{\}]
results of the simulation
    \end{Verbatim}

    \begin{center}
    \adjustimage{max size={0.9\linewidth}{0.9\paperheight}}{trabalho2_files/trabalho2_27_1.png}
    \end{center}
    { \hspace*{\fill} \\}
    
    \begin{tcolorbox}[breakable, size=fbox, boxrule=1pt, pad at break*=1mm,colback=cellbackground, colframe=cellborder]
\prompt{In}{incolor}{167}{\boxspacing}
\begin{Verbatim}[commandchars=\\\{\}]
\PY{n+nb}{print}\PY{p}{(}\PY{l+s+s2}{\PYZdq{}}\PY{l+s+s2}{constant oracle job results:}\PY{l+s+s2}{\PYZdq{}}\PY{p}{)}
\PY{n}{fetch\PYZus{}job\PYZus{}results}\PY{p}{(}\PY{l+s+s2}{\PYZdq{}}\PY{l+s+s2}{cxcrjntztp300083373g}\PY{l+s+s2}{\PYZdq{}}\PY{p}{)}
\end{Verbatim}
\end{tcolorbox}

    \begin{Verbatim}[commandchars=\\\{\}]
constant oracle job results:
    \end{Verbatim}

    \begin{center}
    \adjustimage{max size={0.9\linewidth}{0.9\paperheight}}{trabalho2_files/trabalho2_28_1.png}
    \end{center}
    { \hspace*{\fill} \\}
    
    \hypertarget{grovers-algorithm}{%
\subsubsection{\texorpdfstring{2.5 Grover's Algorithm
}{2.5 Grover's Algorithm }}\label{grovers-algorithm}}

Apresente a sua implementação, simulação de alguns casos, e execução em
hardware real de pelo menos dois casos aproventando o número de qubits
disponiveis.

    \begin{tcolorbox}[breakable, size=fbox, boxrule=1pt, pad at break*=1mm,colback=cellbackground, colframe=cellborder]
\prompt{In}{incolor}{168}{\boxspacing}
\begin{Verbatim}[commandchars=\\\{\}]
\PY{k}{def} \PY{n+nf}{createPhaseOracle\PYZus{}1}\PY{p}{(}\PY{p}{)}\PY{p}{:}
    \PY{c+c1}{\PYZsh{}states: |0101\PYZgt{}, |0110\PYZgt{}, |1100\PYZgt{}, |1111\PYZgt{}}
    \PY{n}{oracle1} \PY{o}{=} \PY{n}{QuantumCircuit}\PY{p}{(}\PY{n+nb}{len}\PY{p}{)}
    \PY{n}{oracle1}\PY{o}{.}\PY{n}{barrier}\PY{p}{(}\PY{p}{)}
    \PY{n}{oracle1}\PY{o}{.}\PY{n}{cz}\PY{p}{(}\PY{l+m+mi}{0}\PY{p}{,} \PY{l+m+mi}{2}\PY{p}{)}
    \PY{n}{oracle1}\PY{o}{.}\PY{n}{cz}\PY{p}{(}\PY{l+m+mi}{1}\PY{p}{,} \PY{l+m+mi}{2}\PY{p}{)}
    \PY{n}{oracle1}\PY{o}{.}\PY{n}{cz}\PY{p}{(}\PY{l+m+mi}{2}\PY{p}{,} \PY{l+m+mi}{3}\PY{p}{)}
    \PY{n}{oracle1}\PY{o}{.}\PY{n}{barrier}\PY{p}{(}\PY{p}{)}
    \PY{k}{return} \PY{n}{oracle1}

\PY{k}{def} \PY{n+nf}{createPhaseOracle\PYZus{}2}\PY{p}{(}\PY{p}{)}\PY{p}{:}
    \PY{c+c1}{\PYZsh{}states: |0110\PYZgt{}, |0111\PYZgt{}, |1100\PYZgt{}, |1101\PYZgt{}}
    \PY{n}{oracle2} \PY{o}{=} \PY{n}{QuantumCircuit}\PY{p}{(}\PY{n+nb}{len}\PY{p}{)}
    \PY{n}{oracle2}\PY{o}{.}\PY{n}{barrier}\PY{p}{(}\PY{p}{)}
    \PY{n}{oracle2}\PY{o}{.}\PY{n}{cz}\PY{p}{(}\PY{l+m+mi}{1}\PY{p}{,} \PY{l+m+mi}{2}\PY{p}{)}
    \PY{n}{oracle2}\PY{o}{.}\PY{n}{cz}\PY{p}{(}\PY{l+m+mi}{2}\PY{p}{,} \PY{l+m+mi}{3}\PY{p}{)}
    \PY{n}{oracle2}\PY{o}{.}\PY{n}{barrier}\PY{p}{(}\PY{p}{)}    
    \PY{k}{return} \PY{n}{oracle2}
\end{Verbatim}
\end{tcolorbox}

    \begin{tcolorbox}[breakable, size=fbox, boxrule=1pt, pad at break*=1mm,colback=cellbackground, colframe=cellborder]
\prompt{In}{incolor}{169}{\boxspacing}
\begin{Verbatim}[commandchars=\\\{\}]
\PY{k}{def} \PY{n+nf}{createDiffuser}\PY{p}{(}\PY{p}{)}\PY{p}{:}
    \PY{n}{diffuser} \PY{o}{=} \PY{n}{QuantumCircuit}\PY{p}{(}\PY{n+nb}{len}\PY{p}{)}
    \PY{n}{diffuser}\PY{o}{.}\PY{n}{barrier}\PY{p}{(}\PY{p}{)}
    
    \PY{k}{for} \PY{n}{qubit} \PY{o+ow}{in} \PY{n+nb}{range}\PY{p}{(}\PY{n+nb}{len}\PY{p}{)}\PY{p}{:}
        \PY{n}{diffuser}\PY{o}{.}\PY{n}{h}\PY{p}{(}\PY{n}{qubit}\PY{p}{)}
        
    \PY{k}{for} \PY{n}{qubit} \PY{o+ow}{in} \PY{n+nb}{range}\PY{p}{(}\PY{n+nb}{len}\PY{p}{)}\PY{p}{:}
        \PY{n}{diffuser}\PY{o}{.}\PY{n}{x}\PY{p}{(}\PY{n}{qubit}\PY{p}{)}
        
    \PY{n}{diffuser}\PY{o}{.}\PY{n}{h}\PY{p}{(}\PY{n+nb}{len}\PY{o}{\PYZhy{}}\PY{l+m+mi}{1}\PY{p}{)}
    \PY{n}{diffuser}\PY{o}{.}\PY{n}{mcx}\PY{p}{(}\PY{n+nb}{list}\PY{p}{(}\PY{n+nb}{range}\PY{p}{(}\PY{n+nb}{len}\PY{o}{\PYZhy{}}\PY{l+m+mi}{1}\PY{p}{)}\PY{p}{)}\PY{p}{,} \PY{n+nb}{len}\PY{o}{\PYZhy{}}\PY{l+m+mi}{1}\PY{p}{)}
    \PY{n}{diffuser}\PY{o}{.}\PY{n}{h}\PY{p}{(}\PY{n+nb}{len}\PY{o}{\PYZhy{}}\PY{l+m+mi}{1}\PY{p}{)}
    
    \PY{k}{for} \PY{n}{qubit} \PY{o+ow}{in} \PY{n+nb}{range}\PY{p}{(}\PY{n+nb}{len}\PY{p}{)}\PY{p}{:}
        \PY{n}{diffuser}\PY{o}{.}\PY{n}{x}\PY{p}{(}\PY{n}{qubit}\PY{p}{)}
        
    \PY{k}{for} \PY{n}{qubit} \PY{o+ow}{in} \PY{n+nb}{range}\PY{p}{(}\PY{n+nb}{len}\PY{p}{)}\PY{p}{:}
        \PY{n}{diffuser}\PY{o}{.}\PY{n}{h}\PY{p}{(}\PY{n}{qubit}\PY{p}{)}
    \PY{n}{diffuser}\PY{o}{.}\PY{n}{barrier}\PY{p}{(}\PY{p}{)}
        
    \PY{k}{return} \PY{n}{diffuser}
\end{Verbatim}
\end{tcolorbox}

    \begin{tcolorbox}[breakable, size=fbox, boxrule=1pt, pad at break*=1mm,colback=cellbackground, colframe=cellborder]
\prompt{In}{incolor}{170}{\boxspacing}
\begin{Verbatim}[commandchars=\\\{\}]
\PY{k}{def} \PY{n+nf}{createCircuit1\PYZus{}grover}\PY{p}{(}\PY{p}{)}\PY{p}{:}
    \PY{n}{grover1} \PY{o}{=} \PY{n}{QuantumCircuit}\PY{p}{(}\PY{n+nb}{len}\PY{p}{)}
    
    \PY{k}{for} \PY{n}{qubit} \PY{o+ow}{in} \PY{n+nb}{range}\PY{p}{(}\PY{n+nb}{len}\PY{p}{)}\PY{p}{:}
        \PY{n}{grover1}\PY{o}{.}\PY{n}{h}\PY{p}{(}\PY{n}{qubit}\PY{p}{)}
    
    \PY{c+c1}{\PYZsh{}target states}
    \PY{n}{grover1} \PY{o}{=} \PY{n}{grover1}\PY{o}{.}\PY{n}{compose}\PY{p}{(}\PY{n}{createPhaseOracle\PYZus{}1}\PY{p}{(}\PY{p}{)}\PY{p}{,} \PY{n+nb}{list}\PY{p}{(}\PY{n+nb}{range}\PY{p}{(}\PY{n+nb}{len}\PY{p}{)}\PY{p}{)}\PY{p}{)}

    \PY{n}{grover1} \PY{o}{=} \PY{n}{grover1}\PY{o}{.}\PY{n}{compose}\PY{p}{(}\PY{n}{createDiffuser}\PY{p}{(}\PY{p}{)}\PY{p}{,} \PY{n+nb}{list}\PY{p}{(}\PY{n+nb}{range}\PY{p}{(}\PY{n+nb}{len}\PY{p}{)}\PY{p}{)}\PY{p}{)}
    \PY{n}{grover1}\PY{o}{.}\PY{n}{measure\PYZus{}all}\PY{p}{(}\PY{p}{)}
    
    \PY{k}{return} \PY{n}{grover1}

\PY{k}{def} \PY{n+nf}{createCircuit2\PYZus{}grover}\PY{p}{(}\PY{p}{)}\PY{p}{:}
    \PY{n}{grover2} \PY{o}{=} \PY{n}{QuantumCircuit}\PY{p}{(}\PY{n+nb}{len}\PY{p}{)}
    
    \PY{k}{for} \PY{n}{qubit} \PY{o+ow}{in} \PY{n+nb}{range}\PY{p}{(}\PY{n+nb}{len}\PY{p}{)}\PY{p}{:}
        \PY{n}{grover2}\PY{o}{.}\PY{n}{h}\PY{p}{(}\PY{n}{qubit}\PY{p}{)}
    
    \PY{c+c1}{\PYZsh{}target states}
    \PY{n}{grover2} \PY{o}{=} \PY{n}{grover2}\PY{o}{.}\PY{n}{compose}\PY{p}{(}\PY{n}{createPhaseOracle\PYZus{}2}\PY{p}{(}\PY{p}{)}\PY{p}{,} \PY{n+nb}{list}\PY{p}{(}\PY{n+nb}{range}\PY{p}{(}\PY{n+nb}{len}\PY{p}{)}\PY{p}{)}\PY{p}{)}
        
    \PY{n}{grover2} \PY{o}{=} \PY{n}{grover2}\PY{o}{.}\PY{n}{compose}\PY{p}{(}\PY{n}{createDiffuser}\PY{p}{(}\PY{p}{)}\PY{p}{,} \PY{n+nb}{list}\PY{p}{(}\PY{n+nb}{range}\PY{p}{(}\PY{n+nb}{len}\PY{p}{)}\PY{p}{)}\PY{p}{)}
    \PY{n}{grover2}\PY{o}{.}\PY{n}{measure\PYZus{}all}\PY{p}{(}\PY{p}{)}
    
    \PY{k}{return} \PY{n}{grover2}
\end{Verbatim}
\end{tcolorbox}

    \begin{tcolorbox}[breakable, size=fbox, boxrule=1pt, pad at break*=1mm,colback=cellbackground, colframe=cellborder]
\prompt{In}{incolor}{171}{\boxspacing}
\begin{Verbatim}[commandchars=\\\{\}]
\PY{k}{def} \PY{n+nf}{runCircuit\PYZus{}grover}\PY{p}{(}\PY{n}{grover}\PY{p}{)}\PY{p}{:}
    \PY{n}{transpiledGrover} \PY{o}{=} \PY{n}{transpile}\PY{p}{(}\PY{n}{grover}\PY{p}{,} \PY{n}{backend}\PY{p}{)}
    \PY{n}{transpiledGrover}\PY{o}{.}\PY{n}{draw}\PY{p}{(}\PY{l+s+s1}{\PYZsq{}}\PY{l+s+s1}{mpl}\PY{l+s+s1}{\PYZsq{}}\PY{p}{,} \PY{n}{idle\PYZus{}wires}\PY{o}{=}\PY{k+kc}{False}\PY{p}{)}
    
    \PY{n}{job} \PY{o}{=} \PY{n}{backend}\PY{o}{.}\PY{n}{run}\PY{p}{(}\PY{p}{[}\PY{n}{transpiledGrover}\PY{p}{]}\PY{p}{,} \PY{n}{shots}\PY{o}{=}\PY{l+m+mi}{1024}\PY{p}{)}
    \PY{n}{job}\PY{o}{.}\PY{n}{job\PYZus{}id}\PY{p}{(}\PY{p}{)}
    
    \PY{n}{results} \PY{o}{=} \PY{n}{job}\PY{o}{.}\PY{n}{result}\PY{p}{(}\PY{p}{)}
    \PY{n}{graph} \PY{o}{=} \PY{n}{results}\PY{o}{.}\PY{n}{get\PYZus{}counts}\PY{p}{(}\PY{p}{)}
    \PY{n}{plot\PYZus{}histogram}\PY{p}{(}\PY{n}{graph}\PY{p}{)}
\end{Verbatim}
\end{tcolorbox}

    \begin{tcolorbox}[breakable, size=fbox, boxrule=1pt, pad at break*=1mm,colback=cellbackground, colframe=cellborder]
\prompt{In}{incolor}{172}{\boxspacing}
\begin{Verbatim}[commandchars=\\\{\}]
\PY{k}{def} \PY{n+nf}{simulateCircuit\PYZus{}grover}\PY{p}{(}\PY{n}{grover}\PY{p}{)}\PY{p}{:}
    \PY{n}{job} \PY{o}{=} \PY{n}{sim}\PY{o}{.}\PY{n}{run}\PY{p}{(}\PY{n}{grover}\PY{p}{)}
    \PY{n}{result} \PY{o}{=} \PY{n}{job}\PY{o}{.}\PY{n}{result}\PY{p}{(}\PY{p}{)}
    \PY{n}{graph} \PY{o}{=} \PY{n}{result}\PY{o}{.}\PY{n}{get\PYZus{}counts}\PY{p}{(}\PY{p}{)}
    \PY{n}{plot\PYZus{}histogram}\PY{p}{(}\PY{n}{graph}\PY{p}{)}
\end{Verbatim}
\end{tcolorbox}

    \begin{tcolorbox}[breakable, size=fbox, boxrule=1pt, pad at break*=1mm,colback=cellbackground, colframe=cellborder]
\prompt{In}{incolor}{173}{\boxspacing}
\begin{Verbatim}[commandchars=\\\{\}]
\PY{c+c1}{\PYZsh{}create grover circuit (with oracle 1)}
\PY{n}{grover\PYZus{}1} \PY{o}{=} \PY{n}{createCircuit1\PYZus{}grover}\PY{p}{(}\PY{p}{)}

\PY{c+c1}{\PYZsh{}run it on real hardware}
\PY{c+c1}{\PYZsh{}runCircuit\PYZus{}grover(grover\PYZus{}1)}

\PY{n+nb}{print}\PY{p}{(}\PY{l+s+s2}{\PYZdq{}}\PY{l+s+s2}{grover\PYZus{}1 circuit:}\PY{l+s+s2}{\PYZdq{}}\PY{p}{)}
\PY{n}{grover\PYZus{}1}\PY{o}{.}\PY{n}{draw}\PY{p}{(}\PY{l+s+s1}{\PYZsq{}}\PY{l+s+s1}{mpl}\PY{l+s+s1}{\PYZsq{}}\PY{p}{,} \PY{n}{idle\PYZus{}wires}\PY{o}{=}\PY{k+kc}{False}\PY{p}{)}
\end{Verbatim}
\end{tcolorbox}

    \begin{Verbatim}[commandchars=\\\{\}]
grover\_1 circuit:
    \end{Verbatim}
 
            
\prompt{Out}{outcolor}{173}{}
    
    \begin{center}
    \adjustimage{max size={0.9\linewidth}{0.9\paperheight}}{trabalho2_files/trabalho2_35_1.png}
    \end{center}
    { \hspace*{\fill} \\}
    

    \begin{center}
    \adjustimage{max size={0.9\linewidth}{0.9\paperheight}}{trabalho2_files/trabalho2_35_2.png}
    \end{center}
    { \hspace*{\fill} \\}
    
    \begin{tcolorbox}[breakable, size=fbox, boxrule=1pt, pad at break*=1mm,colback=cellbackground, colframe=cellborder]
\prompt{In}{incolor}{174}{\boxspacing}
\begin{Verbatim}[commandchars=\\\{\}]
\PY{n+nb}{print}\PY{p}{(}\PY{l+s+s2}{\PYZdq{}}\PY{l+s+s2}{grover\PYZus{}1 simulation:}\PY{l+s+s2}{\PYZdq{}}\PY{p}{)}
\PY{n+nb}{print} \PY{p}{(}\PY{n}{grover\PYZus{}1}\PY{p}{)}
\end{Verbatim}
\end{tcolorbox}

    \begin{Verbatim}[commandchars=\\\{\}]
grover\_1 simulation:
        ┌───┐ ░           ░  ░ ┌───┐┌───┐          ┌───┐┌───┐      ░  ░ ┌─┐   »
   q\_0: ┤ H ├─░──■────────░──░─┤ H ├┤ X ├───────■──┤ X ├┤ H ├──────░──░─┤M├───»
        ├───┤ ░  │        ░  ░ ├───┤├───┤       │  ├───┤├───┤      ░  ░ └╥┘┌─┐»
   q\_1: ┤ H ├─░──┼──■─────░──░─┤ H ├┤ X ├───────■──┤ X ├┤ H ├──────░──░──╫─┤M├»
        ├───┤ ░  │  │     ░  ░ ├───┤├───┤       │  ├───┤├───┤      ░  ░  ║ └╥┘»
   q\_2: ┤ H ├─░──■──■──■──░──░─┤ H ├┤ X ├───────■──┤ X ├┤ H ├──────░──░──╫──╫─»
        ├───┤ ░        │  ░  ░ ├───┤├───┤┌───┐┌─┴─┐├───┤├───┤┌───┐ ░  ░  ║  ║ »
   q\_3: ┤ H ├─░────────■──░──░─┤ H ├┤ X ├┤ H ├┤ X ├┤ H ├┤ X ├┤ H ├─░──░──╫──╫─»
        └───┘ ░           ░  ░ └───┘└───┘└───┘└───┘└───┘└───┘└───┘ ░  ░  ║  ║ »
meas: 4/═════════════════════════════════════════════════════════════════╩══╩═»
                                                                         0  1 »
«
«   q\_0: ──────
«
«   q\_1: ──────
«        ┌─┐
«   q\_2: ┤M├───
«        └╥┘┌─┐
«   q\_3: ─╫─┤M├
«         ║ └╥┘
«meas: 4/═╩══╩═
«         2  3
    \end{Verbatim}

    \begin{tcolorbox}[breakable, size=fbox, boxrule=1pt, pad at break*=1mm,colback=cellbackground, colframe=cellborder]
\prompt{In}{incolor}{175}{\boxspacing}
\begin{Verbatim}[commandchars=\\\{\}]
\PY{n+nb}{print}\PY{p}{(}\PY{l+s+s2}{\PYZdq{}}\PY{l+s+s2}{grover\PYZus{}1 simulation results:}\PY{l+s+s2}{\PYZdq{}}\PY{p}{)}
\PY{n}{simulateCircuit\PYZus{}grover}\PY{p}{(}\PY{n}{grover\PYZus{}1}\PY{p}{)}
\end{Verbatim}
\end{tcolorbox}

    \begin{Verbatim}[commandchars=\\\{\}]
grover\_1 simulation results:
    \end{Verbatim}

    \begin{center}
    \adjustimage{max size={0.9\linewidth}{0.9\paperheight}}{trabalho2_files/trabalho2_37_1.png}
    \end{center}
    { \hspace*{\fill} \\}
    
    \begin{tcolorbox}[breakable, size=fbox, boxrule=1pt, pad at break*=1mm,colback=cellbackground, colframe=cellborder]
\prompt{In}{incolor}{176}{\boxspacing}
\begin{Verbatim}[commandchars=\\\{\}]
\PY{n+nb}{print}\PY{p}{(}\PY{l+s+s2}{\PYZdq{}}\PY{l+s+s2}{grover\PYZus{}1 job results:}\PY{l+s+s2}{\PYZdq{}}\PY{p}{)}
\PY{n}{fetch\PYZus{}job\PYZus{}results}\PY{p}{(}\PY{l+s+s2}{\PYZdq{}}\PY{l+s+s2}{cxcrjtvztp3000833740}\PY{l+s+s2}{\PYZdq{}}\PY{p}{)}
\end{Verbatim}
\end{tcolorbox}

    \begin{Verbatim}[commandchars=\\\{\}]
grover\_1 job results:
    \end{Verbatim}

    \begin{center}
    \adjustimage{max size={0.9\linewidth}{0.9\paperheight}}{trabalho2_files/trabalho2_38_1.png}
    \end{center}
    { \hspace*{\fill} \\}
    
    \begin{tcolorbox}[breakable, size=fbox, boxrule=1pt, pad at break*=1mm,colback=cellbackground, colframe=cellborder]
\prompt{In}{incolor}{177}{\boxspacing}
\begin{Verbatim}[commandchars=\\\{\}]
\PY{c+c1}{\PYZsh{} create grover circuit (with oracle 2)}
\PY{n}{grover\PYZus{}2} \PY{o}{=} \PY{n}{createCircuit2\PYZus{}grover}\PY{p}{(}\PY{p}{)}

\PY{c+c1}{\PYZsh{}run it on real hardware}
\PY{c+c1}{\PYZsh{}runCircuit\PYZus{}grover(grover\PYZus{}2)}

\PY{n+nb}{print}\PY{p}{(}\PY{l+s+s2}{\PYZdq{}}\PY{l+s+s2}{grover\PYZus{}2 circuit:}\PY{l+s+s2}{\PYZdq{}}\PY{p}{)}
\PY{n}{grover\PYZus{}2}\PY{o}{.}\PY{n}{draw}\PY{p}{(}\PY{l+s+s1}{\PYZsq{}}\PY{l+s+s1}{mpl}\PY{l+s+s1}{\PYZsq{}}\PY{p}{,} \PY{n}{idle\PYZus{}wires}\PY{o}{=}\PY{k+kc}{False}\PY{p}{)}
\end{Verbatim}
\end{tcolorbox}

    \begin{Verbatim}[commandchars=\\\{\}]
grover\_2 circuit:
    \end{Verbatim}
 
            
\prompt{Out}{outcolor}{177}{}
    
    \begin{center}
    \adjustimage{max size={0.9\linewidth}{0.9\paperheight}}{trabalho2_files/trabalho2_39_1.png}
    \end{center}
    { \hspace*{\fill} \\}
    

    \begin{center}
    \adjustimage{max size={0.9\linewidth}{0.9\paperheight}}{trabalho2_files/trabalho2_39_2.png}
    \end{center}
    { \hspace*{\fill} \\}
    
    \begin{tcolorbox}[breakable, size=fbox, boxrule=1pt, pad at break*=1mm,colback=cellbackground, colframe=cellborder]
\prompt{In}{incolor}{178}{\boxspacing}
\begin{Verbatim}[commandchars=\\\{\}]
\PY{n+nb}{print}\PY{p}{(}\PY{l+s+s2}{\PYZdq{}}\PY{l+s+s2}{grover\PYZus{}2 simulation:}\PY{l+s+s2}{\PYZdq{}}\PY{p}{)}
\PY{n+nb}{print} \PY{p}{(}\PY{n}{grover\PYZus{}2}\PY{p}{)}
\end{Verbatim}
\end{tcolorbox}

    \begin{Verbatim}[commandchars=\\\{\}]
grover\_2 simulation:
        ┌───┐ ░        ░  ░ ┌───┐┌───┐          ┌───┐┌───┐      ░  ░ ┌─┐      »
   q\_0: ┤ H ├─░────────░──░─┤ H ├┤ X ├───────■──┤ X ├┤ H ├──────░──░─┤M├──────»
        ├───┤ ░        ░  ░ ├───┤├───┤       │  ├───┤├───┤      ░  ░ └╥┘┌─┐   »
   q\_1: ┤ H ├─░──■─────░──░─┤ H ├┤ X ├───────■──┤ X ├┤ H ├──────░──░──╫─┤M├───»
        ├───┤ ░  │     ░  ░ ├───┤├───┤       │  ├───┤├───┤      ░  ░  ║ └╥┘┌─┐»
   q\_2: ┤ H ├─░──■──■──░──░─┤ H ├┤ X ├───────■──┤ X ├┤ H ├──────░──░──╫──╫─┤M├»
        ├───┤ ░     │  ░  ░ ├───┤├───┤┌───┐┌─┴─┐├───┤├───┤┌───┐ ░  ░  ║  ║ └╥┘»
   q\_3: ┤ H ├─░─────■──░──░─┤ H ├┤ X ├┤ H ├┤ X ├┤ H ├┤ X ├┤ H ├─░──░──╫──╫──╫─»
        └───┘ ░        ░  ░ └───┘└───┘└───┘└───┘└───┘└───┘└───┘ ░  ░  ║  ║  ║ »
meas: 4/══════════════════════════════════════════════════════════════╩══╩══╩═»
                                                                      0  1  2 »
«
«   q\_0: ───
«
«   q\_1: ───
«
«   q\_2: ───
«        ┌─┐
«   q\_3: ┤M├
«        └╥┘
«meas: 4/═╩═
«         3
    \end{Verbatim}

    \begin{tcolorbox}[breakable, size=fbox, boxrule=1pt, pad at break*=1mm,colback=cellbackground, colframe=cellborder]
\prompt{In}{incolor}{179}{\boxspacing}
\begin{Verbatim}[commandchars=\\\{\}]
\PY{n+nb}{print}\PY{p}{(}\PY{l+s+s2}{\PYZdq{}}\PY{l+s+s2}{grover\PYZus{}2 simulation results:}\PY{l+s+s2}{\PYZdq{}}\PY{p}{)}
\PY{n}{simulateCircuit\PYZus{}grover}\PY{p}{(}\PY{n}{grover\PYZus{}2}\PY{p}{)}
\end{Verbatim}
\end{tcolorbox}

    \begin{Verbatim}[commandchars=\\\{\}]
grover\_2 simulation results:
    \end{Verbatim}

    \begin{center}
    \adjustimage{max size={0.9\linewidth}{0.9\paperheight}}{trabalho2_files/trabalho2_41_1.png}
    \end{center}
    { \hspace*{\fill} \\}
    
    \begin{tcolorbox}[breakable, size=fbox, boxrule=1pt, pad at break*=1mm,colback=cellbackground, colframe=cellborder]
\prompt{In}{incolor}{180}{\boxspacing}
\begin{Verbatim}[commandchars=\\\{\}]
\PY{n+nb}{print}\PY{p}{(}\PY{l+s+s2}{\PYZdq{}}\PY{l+s+s2}{grover\PYZus{}2 job results:}\PY{l+s+s2}{\PYZdq{}}\PY{p}{)}
\PY{n}{fetch\PYZus{}job\PYZus{}results}\PY{p}{(}\PY{l+s+s2}{\PYZdq{}}\PY{l+s+s2}{cxcrk0ctpsjg008mk13g}\PY{l+s+s2}{\PYZdq{}}\PY{p}{)}
\end{Verbatim}
\end{tcolorbox}

    \begin{Verbatim}[commandchars=\\\{\}]
grover\_2 job results:
    \end{Verbatim}

    \begin{center}
    \adjustimage{max size={0.9\linewidth}{0.9\paperheight}}{trabalho2_files/trabalho2_42_1.png}
    \end{center}
    { \hspace*{\fill} \\}
    

    % Add a bibliography block to the postdoc
    
    
    
\end{document}
